\begin{thebibliography}{99}\setlength{\itemsep}{2pt}\small
\bibitem[21Shares(2025)]{shares2025} 21Shares (2025). Record crypto liquidations amid tariff shock. Research note, October 13, 2025. \url{https://www.21shares.com/research/record-crypto-liquidations-amid-tariff-shock}.
\bibitem[Achieser(1956)]{achieser1956} Achieser, N. I. (1956). \emph{Theory of Approximation}. New York: Frederick Ungar.
\bibitem[Almgren and Chriss(2001)]{almgren2001} Almgren, R., and N. Chriss (2001). Optimal execution of portfolio transactions. \emph{Journal of Risk} 3, 5--39.
\bibitem[Amihud(2002)]{amihud2002} Amihud, Y. (2002). Illiquidity and stock returns: Cross-section and time-series effects. \emph{Journal of Financial Markets} 5, 31--56.
\bibitem[Andersen and Bondarenko(2014)]{andersen2014} Andersen, T. G., and O. Bondarenko (2014). VPIN and the flash crash. \emph{Journal of Financial Markets} 17, 1--46.
\bibitem[Avellaneda and Stoikov(2008)]{avellaneda2008} Avellaneda, M., and S. Stoikov (2008). High-frequency trading in a limit order book. \emph{Quantitative Finance} 8, 217--224.
\bibitem[Biais et~al.(2005)]{biais2005} Biais, B., L. Glosten, and C. Spatt (2005). Market microstructure: A survey of microfoundations, empirical results, and policy implications. \emph{Journal of Financial Markets} 8, 217--264.
\bibitem[Biais et~al.(1995)]{biais1995} Biais, B., P. Hillion, and C. Spatt (1995). An empirical analysis of the limit order book and the order flow in the Paris Bourse. \emph{Journal of Finance} 50, 1655--1689.
\bibitem[Bouchaud et~al.(2004)]{bouchaud2004} Bouchaud, J.-P., Y. Gefen, M. Potters, and M. Wyart (2004). Fluctuations and response in financial markets: The subtle nature of ``random'' price changes. \emph{Quantitative Finance} 4, 176--190.
\bibitem[Bouchaud et~al.(2018)]{bouchaud2018} Bouchaud, J.-P., J. Bonart, J. Donier, and M. Gould (2018). \emph{Trades, Quotes and Prices: Financial Markets Under the Microscope}. Cambridge University Press.
\bibitem[Brogaard et~al.(2014)]{brogaard2014} Brogaard, J., T. Hendershott, and R. Riordan (2014). High-frequency trading and price discovery. \emph{Review of Financial Studies} 27, 2267--2306.
\bibitem[Brunnermeier and Pedersen(2009)]{brunnermeier2009} Brunnermeier, M. K., and L. H. Pedersen (2009). Market liquidity and funding liquidity. \emph{Review of Financial Studies} 22, 2201--2238.
\bibitem[Budish et~al.(2015)]{budish2015} Budish, E., P. Cramton, and J. Shim (2015). The high-frequency trading arms race: Frequent batch auctions as a market design response. \emph{Quarterly Journal of Economics} 130, 1547--1621.
\bibitem[Cartea et~al.(2015)]{cartea2015} Cartea, \'A., S. Jaimungal, and J. Penalva (2015). \emph{Algorithmic and High-Frequency Trading}. Cambridge University Press.
\bibitem[Cespa and Foucault(2014)]{cespa2014} Cespa, G., and T. Foucault (2014). Illiquidity contagion and liquidity crashes. \emph{Review of Financial Studies} 27, 1615--1660.
\bibitem[Clark and West(2007)]{clark2007} Clark, T. E., and K. D. West (2007). Approximately normal tests for equal predictive accuracy in nested models. \emph{Journal of Econometrics} 138, 291--311.
\bibitem[Conley and Taber(2011)]{conley2011} Conley, T. G., and C. R. Taber (2011). Inference with ``difference in differences'' with a small number of policy changes. \emph{Review of Economics and Statistics} 93, 113--125.
\bibitem[Cont et~al.(2014)]{cont2014} Cont, R., A. Kukanov, and S. Stoikov (2014). The price impact of order book events. \emph{Journal of Financial Econometrics} 12, 47--88.
\bibitem[Corsi(2009)]{corsi2009} Corsi, F. (2009). A simple approximate long-memory model of realized volatility. \emph{Journal of Financial Econometrics} 7, 174--196.
\bibitem[Dahleh and Diaz-Bobillo(1995)]{dahleh1995} Dahleh, M. A., and I. J. Diaz-Bobillo (1995). \emph{Control of Uncertain Systems: A Linear Programming Approach}. Prentice Hall.
\bibitem[Dahleh and Pearson(1987)]{dahleh1987} Dahleh, M. A., and J. B. Pearson (1987). $\ell^1$-optimal feedback controllers for MIMO discrete-time systems. \emph{IEEE Transactions on Automatic Control} 32, 314--322.
\bibitem[Degryse et~al.(2005)]{degryse2005} Degryse, H., F. de Jong, M. van Ravenswaaij, and G. Wuyts (2005). Aggressive orders and the resiliency of a limit order market. \emph{Review of Finance} 9, 201--242.
\bibitem[Diebold and Mariano(1995)]{diebold1995} Diebold, F. X., and R. S. Mariano (1995). Comparing predictive accuracy. \emph{Journal of Business \& Economic Statistics} 13, 253--263.
\bibitem[Donoho(1994)]{donoho1994} Donoho, D.~L. (1994). Statistical estimation and optimal recovery. \emph{Annals of Statistics} 22, 238--270.
\bibitem[Easley et~al.(2012)]{easley2012} Easley, D., M. M. L\'opez de Prado, and M. O'Hara (2012). Flow toxicity and liquidity in a high-frequency world. \emph{Review of Financial Studies} 25, 1457--1493.
\bibitem[Foucault et~al.(2005)]{foucault2005} Foucault, T., O. Kadan, and E. Kandel (2005). Limit order book as a market for liquidity. \emph{Review of Financial Studies} 18, 1171--1217.
\bibitem[Hansen and Walster(2004)]{hansen2004} Hansen, E., and G. W. Walster (2004). \emph{Global Optimization Using Interval Analysis}, 2nd ed. New York: Marcel Dekker.
\bibitem[Harvey(2017)]{harvey2017} Harvey, C. R. (2017). Presidential address: The scientific outlook in financial economics. \emph{Journal of Finance} 72, 1399--1440.
\bibitem[Harvey et~al.(2016)]{harvey2016} Harvey, C. R., Y. Liu, and H. Zhu (2016). \ldots and the cross-section of expected returns. \emph{Review of Financial Studies} 29, 5--68.
\bibitem[Hendershott et~al.(2011)]{hendershott2011} Hendershott, T., C. M. Jones, and A. J. Menkveld (2011). Does algorithmic trading improve liquidity? \emph{Journal of Finance} 66, 1--33.
\bibitem[Ho and Stoll(1981)]{ho1981} Ho, T., and H. R. Stoll (1981). Optimal dealer pricing under transactions and return uncertainty. \emph{Journal of Financial Economics} 9, 47--73.
\bibitem[Holm(1979)]{holm1979} Holm, S. (1979). A simple sequentially rejective multiple test procedure. \emph{Scandinavian Journal of Statistics} 6, 65--70.
\bibitem[Hotelling(1931)]{hotelling1931} Hotelling, H. (1931). The generalization of Student's ratio. \emph{Annals of Mathematical Statistics} 2, 360--378.
\bibitem[Howard et~al.(2021)]{howard2021} Howard, S. R., A. Ramdas, J. McAuliffe, and J. Sekhon (2021). Time-uniform, nonparametric, nonasymptotic confidence sequences. \emph{Annals of Statistics} 49, 1055--1080.
\bibitem[Huang and Wang(2009)]{huang2009} Huang, J., and J. Wang (2009). Liquidity and market crashes. \emph{Review of Financial Studies} 22, 2607--2643.
\bibitem[Jord\`a(2005)]{jorda2005} Jord\`a, \`O. (2005). Estimation and inference of impulse responses by local projections. \emph{American Economic Review} 95, 161--182.
\bibitem[J.P. Morgan/Reuters(1996)]{riskmetrics1996} J.P. Morgan/Reuters (1996). \emph{RiskMetrics---Technical Document}, 4th ed. New York: J.P. Morgan.
\bibitem[Karlin and Studden(1966)]{karlin1966} Karlin, S., and W. J. Studden (1966). \emph{Tchebycheff Systems: With Applications in Analysis and Statistics}. New York: Interscience.
\bibitem[Kirilenko et~al.(2017)]{kirilenko2017} Kirilenko, A., A. S. Kyle, M. Samadi, and T. Tuzun (2017). The Flash Crash: High-frequency trading in an electronic market. \emph{Journal of Finance} 72, 967--998.
\bibitem[Krein and Nudel'man(1977)]{krein1977} Krein, M. G., and A. A. Nudel'man (1977). \emph{The Markov Moment Problem and Extremal Problems}. Translations of Mathematical Monographs 50. Providence, RI: American Mathematical Society.
\bibitem[Kreiss(1962)]{kreiss1962} Kreiss, H.-O. (1962). \"Uber die Stabilit\"atsdefinition f\"ur Differenzengleichungen die partielle Differentialgleichungen approximieren. \emph{BIT Numerical Mathematics} 2, 153--181.
\bibitem[K\"unsch(1989)]{kunsch1989} K\"unsch, H. R. (1989). The jackknife and the bootstrap for general stationary observations. \emph{Annals of Statistics} 17, 1217--1241.
\bibitem[Kyle(1985)]{kyle1985} Kyle, A. S. (1985). Continuous auctions and insider trading. \emph{Econometrica} 53, 1315--1335.
\bibitem[Kyle and Obizhaeva(2016)]{kyle2016} Kyle, A. S., and A. A. Obizhaeva (2016). Market microstructure invariance: Empirical hypotheses. \emph{Econometrica} 84, 1345--1404.
\bibitem[Large(2007)]{large2007} Large, J. (2007). Measuring the resiliency of an electronic limit order book. \emph{Journal of Financial Markets} 10, 1--25.
\bibitem[Le Cam(1986)]{lecam1986} Le Cam, L. (1986). \emph{Asymptotic Methods in Statistical Decision Theory}. New York: Springer.
\bibitem[Makarov and Schoar(2020)]{makarov2020} Makarov, I., and A. Schoar (2020). Trading and arbitrage in cryptocurrency markets. \emph{Journal of Financial Economics} 135, 293--319.
\bibitem[Menkveld(2013)]{menkveld2013} Menkveld, A. J. (2013). High frequency trading and the new market makers. \emph{Journal of Financial Markets} 16, 712--740.
\bibitem[Milanese and Vicino(1991)]{milanese1991} Milanese, M., and A. Vicino (1991). Optimal estimation theory for dynamic systems with set membership uncertainty: An overview. \emph{Automatica} 27, 997--1009.
\bibitem[Newey and West(1987)]{newey1987} Newey, W. K., and K. D. West (1987). A simple, positive semi-definite, heteroskedasticity and autocorrelation consistent covariance matrix. \emph{Econometrica} 55, 703--708.
\bibitem[Niroula et~al.(2026)]{niroula2026} Niroula, P., P. Niroula, and S. Khatiwada (2026). What prices reveal about hidden trader populations, and what identifying them costs. Working paper, Minnesota State University, Mankato.
\bibitem[Nosek et~al.(2018)]{nosek2018} Nosek, B. A., C. R. Ebersole, A. C. DeHaven, and D. T. Mellor (2018). The preregistration revolution. \emph{Proceedings of the National Academy of Sciences} 115, 2600--2606.
\bibitem[Obizhaeva and Wang(2013)]{obizhaeva2013} Obizhaeva, A. A., and J. Wang (2013). Optimal trading strategy and supply/demand dynamics. \emph{Journal of Financial Markets} 16, 1--32.
\bibitem[Rump(2010)]{rump2010} Rump, S. M. (2010). Verification methods: Rigorous results using floating-point arithmetic. \emph{Acta Numerica} 19, 287--449.
\bibitem[Schneider and Lillo(2019)]{schneider2019} Schneider, M., and F. Lillo (2019). Cross-impact and no-dynamic-arbitrage. \emph{Quantitative Finance} 19, 137--154.
\bibitem[SEC(2010)]{sec2010} U.S. Securities and Exchange Commission (2010). Risk management controls for brokers or dealers with market access. Release No.\ 34-63241.
\bibitem[Shleifer and Vishny(1992)]{shleifer1992} Shleifer, A., and R. W. Vishny (1992). Liquidation values and debt capacity: A market equilibrium approach. \emph{Journal of Finance} 47, 1343--1366.
\bibitem[Sornette et~al.(2023)]{sornette2023} Sornette, D., S. Lera, J. Lin, and K. Wu (2023). Non-normal interactions create socio-economic bubbles. \emph{Communications Physics} 6, 261.
\bibitem[Spijker(1991)]{spijker1991} Spijker, M. N. (1991). On a conjecture by LeVeque and Trefethen related to the Kreiss matrix theorem. \emph{BIT Numerical Mathematics} 31, 551--555.
\bibitem[Trefethen and Embree(2005)]{trefethen2005} Trefethen, L. N., and M. Embree (2005). \emph{Spectra and Pseudospectra: The Behavior of Nonnormal Matrices and Operators}. Princeton University Press.
\bibitem[Tu et~al.(2018)]{tu2018} Tu, S., R. Boczar, and B. Recht (2018). Minimax lower bounds for $\mathcal H_\infty$-norm estimation. Working paper, arXiv:1809.10855.
\bibitem[Vidyasagar(1986)]{vidyasagar1986} Vidyasagar, M. (1986). Optimal rejection of persistent bounded disturbances. \emph{IEEE Transactions on Automatic Control} 31, 527--534.
\bibitem[Ville(1939)]{ville1939} Ville, J. (1939). \emph{\'Etude critique de la notion de collectif}. Paris: Gauthier-Villars.
\bibitem[van Waarde et~al.(2020)]{vanwaarde2020} van Waarde, H. J., J. Eising, H. L. Trentelman, and M. K. Camlibel (2020). Data informativity: A new perspective on data-driven analysis and control. \emph{IEEE Transactions on Automatic Control} 65, 4753--4768.
\end{thebibliography}

\appendix
\section{Proofs}\label{app:proofs}

\paragraph{Proposition~\ref{prop:jury}.} With one cohort, $J=\begin{psmallmatrix}a&\kappa q\\ \eta a&1-\eta+\eta\kappa q\end{psmallmatrix}$, so ${\rm tr}J=a+1-\eta+\eta\kappa q$ and $\det J=a(1-\eta)$. The Jury conditions for a real $2\times2$ matrix are $|\det J|<1$ and $1\pm{\rm tr}J+\det J>0$. Substituting $a=1-\kappa\varphi_f$: $1-{\rm tr}J+\det J=\eta(1-a-\kappa q)=\eta\kappa(\varphi_f-q)$, positive iff $q<\varphi_f$; and $1+{\rm tr}J+\det J=(4-2\eta)-\kappa[(2-\eta)\varphi_f-\eta q]$, positive iff the stated bound on $\kappa$ holds (automatically when $(2-\eta)\varphi_f\le\eta q$).

\paragraph{Proposition~\ref{prop:invariance}.} Equal merged measures give equal $G$, hence equal $H$, and the impulse coefficients are the Laurent coefficients of $H$ (companion paper, Theorem on merged equivalence). Under $z\mapsto Tz$, $c^\top J^kb$ is unchanged because $CT^{-1}(TJT^{-1})^kTb=CJ^kb$, while the row norm $\|CT^{-1}(TJT^{-1})^k\|$ depends on $T$. The displayed price path is the convolution of inputs with $\{g_k\}$, and its worst case over the $\ell_\infty$ ball is $\bar u\sum|g|$ by H\"older.

\paragraph{Proposition~\ref{prop:inventory_S}.} (i) Taking $z$-transforms of $y_{i,t}=(1-\eta_i)y_{i,t-1}+\eta_i x_t$, $I_{i,t}=(1-\delta_{\rm inv})I_{i,t-1}+u_{i,t}$ and $u_{i,t}=q_i y_{i,t}-\psi_i I_{i,t-1}$ gives $Y_i(z)=\frac{\eta_i z}{z-1+\eta_i}X(z)$ and $(z-1+\delta_{\rm inv}+\psi_i)I_i(z)=q_i z Y_i(z)$, whence $U_i(z)=\frac{z-1+\delta_{\rm inv}}{z-1+\delta_{\rm inv}+\psi_i}q_i Y_i(z)$ and $G_{\rm inv}(z)=\sum_i U_i(z)/X(z)$ has the stated form. Evaluating at $z=1$ when the augmented operator is Schur stable gives $G_{\rm inv}(1)=\sum_i \frac{\delta_{\rm inv}}{\delta_{\rm inv}+\psi_i}q_i=\Phi_{\rm eff}(\psi)$ and $S(\psi)=H_{\rm inv}(1)=\kappa/(1-a-\kappa\Phi_{\rm eff}(\psi))$. (ii) Write $r_i=1-\eta_i$ and $\rho_i=1-\delta_{\rm inv}-\psi_i$. When $0\le\psi_i\le 1-\delta_{\rm inv}$ and $\delta_{\rm inv}\ge\eta_i$, we have $0\le\rho_i\le 1-\delta_{\rm inv}\le r_i<1$. For $r_i>\rho_i$, the $m$-th open-loop tap is
$k_{i,m}=q_i\eta_i[r_i^m-\psi_i(r_i^m-\rho_i^m)/(r_i-\rho_i)]
=q_i\eta_i[(1-v_i)r_i^m+v_i\rho_i^m]$, where
$v_i=\psi_i/(r_i-\rho_i)=\psi_i/(\delta_{\rm inv}-\eta_i+\psi_i)\in[0,1]$.
It is nonnegative. If $r_i=\rho_i$, the assumptions force $\psi_i=0$ and the tap is $q_i\eta_i r_i^m\ge0$ (with the usual $r_i^0=1$ convention).
Nonnegative open-loop taps $k_m(\psi)\ge0$ with $a+\kappa\sum_m k_m(\psi)=a+\kappa\Phi_{\rm eff}(\psi)<1$ generate a nonnegative closed-loop renewal sequence $h_t(\psi)\ge0$, so $\|h(\psi)\|_1=\sum_{t\ge0}h_t(\psi)=S(\psi)$.

\paragraph{Theorem~\ref{thm:radius}.} (i) $\sup_t\|CJ^tz_0\|\le\sup_t\|CJ^t\|\|z_0\|$, with equality approached by aligning $z_0$ with the top right singular vector of $CJ^{t^*}$ at the maximizing $t^*$. (ii) The Kreiss matrix theorem with Spijker's sharp constant gives $\mathcal K(J)\le\sup_t\|J^t\|\le en\mathcal K(J)$; substitute in (i). (iii) Take $J_\theta=\begin{psmallmatrix}\rho&\theta\\0&\rho/2\end{psmallmatrix}$; then $\rho(J_\theta)=\rho$ and $(1,0)J_\theta e_2=\theta$, so $M_C\ge|\theta|$ and $r^*_{\rm exact}\le B/|\theta|\to0$.

\paragraph{Proposition~\ref{prop:multi}.} With all blocks entrywise nonnegative, every entry of $J_{\rm multi}^t$ is a polynomial in the entries of $K_{21}$ with nonnegative coefficients, so it is non-decreasing in $K_{21}$; with $C\ge0$, so is every entry of $CJ^t$ and hence its norm. For the Kreiss constant, the Neumann series gives $|(zI-J)^{-1}|\le(|z|I-J)^{-1}$ entrywise for $|z|>1$, so the supremum is attained on the real axis, where $(rI-J)^{-1}=\sum_kJ^k/r^{k+1}$ is entrywise non-decreasing in $K_{21}$.

\paragraph{Theorem~\ref{thm:forward}.} After the step, inventory is at most $|I^{\rm post}_{t+1}(\alpha_t)|$ in absolute value. The price move over the step and the cumulative window is bounded by $M_C\|z_t\|_2$ (drift from the current state) plus $s(\alpha)G_{\rm cum}\bar u_{\rm cum}$ (admitted shocks), by superposition and H\"older. The margin requirement scales the loss by $(1+\mu)$. Hedging cost and the post-fill margin increment are deducted directly. Inequality \eqref{eq:guardineq} therefore implies $B_{{\rm rem},t+1}\ge R_{\rm safe}$ for every admissible disturbance, and induction over steps at which an admissible $\alpha$ exists gives forward solvency. Setting $\alpha=0$ ($M_C(0)=1$, no fills, no hedge) gives $\bar u_{\rm adm}$. Sufficiency over the Euclidean ball follows from the norm bounds. An aligned shock can attain a price-response component at rest, but that does not establish joint attainment with the inventory, fill and cost envelopes. The theorem therefore states sufficiency only.

\paragraph{Proposition~\ref{prop:driftguard}.} (i) Positivity gives $|c^\top J^tz_0|\le c^\top J^t\bar z_0$, with equality for memories aligned with $x_0$ (either sign if $x_0=0$). Summing the original memory recursion, without rescaling its initial condition, gives $\sum_{t\ge0}y_{i,t}=y_{i,0}/\eta_i+\sum_{t\ge1}x_t$. Hence $(1-a-\alpha\kappa\Phi)\sum_{t\ge1}x_t=ax_0+\alpha\kappa\sum_iq_iy_{i,0}/\eta_i$, proving \eqref{eq:driftenv}. (ii) Superposition and positivity give the expression in \eqref{eq:driftpeak}; shocks and memories with the common sign of $x_0$ attain it. Taking the maximum over the same time index proves sharpness for fixed parameters. (iii) If $|x_t|,|y_{i,t}|\le L$, then $|x_{t+1}|\le aL+\alpha\kappa\Phi^+L+\kappa\bar u\le L$. Each subsequent memory is a convex combination of two values bounded by $L$. Induction proves the invariant box. Rearranging gives \eqref{eq:alphabox}. For $\bar u>0$, its endpoint automatically has $a+\alpha\kappa\Phi^+<1$. No general maximality or positive capacity cost is inferred by adding separately maximized cumulative and pointwise quantities.

\paragraph{Proposition~\ref{prop:dominance}.} The eigenvalue rule admits $\alpha_\rho$; the one-step certificate fails at $\alpha_\rho$ exactly when the left side of \eqref{eq:guardineq}, with the shock term replaced by its one-step bound $M_C(\alpha_\rho)s(\alpha_\rho)\bar u_{{\rm step},t}$, exceeds $B_{{\rm rem},t}$, which rearranges to the stated inequality. The summed-gain guard admits only levels satisfying \eqref{eq:guardineq}, so it never admits the breaching level, and by the feasibility assumption it admits at least $\alpha=0$.

\paragraph{Proposition~\ref{prop:lucas}.} (i) The steady-state gain is $\eta=P/(P+\sigma_\nu^2)$ with $P$ the positive root of $P^2-\sigma_\zeta^2P-\sigma_\zeta^2\sigma_\nu^2=0$. Since $P$ is strictly increasing in $\sigma_\zeta^2$ and $\eta$ is strictly increasing in $P$, $\eta^*$ is strictly increasing in $\sigma_\zeta^2(\alpha)=\alpha^2\sigma_{\zeta,0}^2$, hence in $\alpha$. The second statement is immediate under the stated monotonicity.

(ii) Write the recursion as a renewal equation: $h_t/\kappa=u_t$ with $u_0=1$ and $u_t=\sum_{j=1}^tp_ju_{t-j}$, where $p_1=a+\kappa k_0$ and $p_j=\kappa k_{j-1}$ for $j\ge2$. The $p_j$ are nonnegative with total mass $\sum_jp_j=a+\kappa\Phi=a+\kappa\sum_iq_i<1$, which does not depend on the adaptation rates, so $u_t$ is the probability that a renewal process whose (defective) inter-arrival law is $p$ has a renewal at time $t$, and $\sum_{t<H}u_t$ is the expected number of renewals in $\{0,\ldots,H-1\}$. The cumulative law is $F(j)=\sum_{i\le j}p_i=a+\kappa\sum_{m<j}k_m=a+\kappa\sum_iq_i\bigl[1-(1-\eta_i)^j\bigr]$, with $\partial F(j)/\partial\eta_i=\kappa q_ij(1-\eta_i)^{j-1}\ge0$. If $\eta'\ge\eta$ coordinatewise, then $F'\ge F$ pointwise with the same total mass, so the quantile coupling $X=F^{-1}(U)$, $X'=F'^{-1}(U)$ ($X=\infty$ when $U>F(\infty)$) of independent inter-arrival times gives $X'\le X$ and hence, pathwise, at least as many renewals before $H$ under $\eta'$. Taking expectations gives $S_H(\eta')\ge S_H(\eta)$. If $q_i>0$ and $\eta'_i>\eta_i$, then $F'(1)>F(1)$, so with positive probability the first renewal occurs at time 1 under $\eta'$ and later under $\eta$, and the inequality is strict for $H\ge2$. As $H\to\infty$ both sides tend to $\kappa/(1-a-\kappa\Phi)$, which does not depend on $\eta$. Combined with (i), throttling lowers every rate, so $S_H(J(\alpha,\eta^*(\alpha)))\le S_H(J(\alpha,\eta^*(1)))$.

\paragraph{Theorem~\ref{thm:window}.} The recursion $h_{t+1}=ah_t+\kappa\sum_{m\le t}k_mh_{t-m}$ with $h_0=\kappa\neq0$ gives $k_t$ from $h_0,\ldots,h_{t+1}$, so the window determines $k_0,\ldots,k_{n-1}$ and nothing more about $M$: any nonnegative measure on $[0,\bar r]$ with these moments, written as atoms $(r_i,w_i)$ or as a limit of such, is the merged measure of a positive economy with $\eta_i=1-r_i$ and $q_i=w_i/\eta_i$ that produces the same window. (i) The compatible measures form a convex, weak-* compact set, and $\Phi=\int f\,dM$ is linear and continuous on it ($f$ is continuous on $[0,\bar r]$), so its range is a closed interval $[\Phi_-,\Phi_+]$. Since $f^{(n)}(r)=n!/(1-r)^{n+1}>0$ on $[0,\bar r]$, $(1,r,\ldots,r^{n-1},f)$ is an extended complete Chebyshev system, and by the Markov--Krein theorem \citep[Theorems IV.1.1 and IV.6.1]{krein1977} the maximum is attained at the upper and the minimum at the lower principal representation, which have at most $\lfloor n/2\rfloor+1$ atoms. The theorem applies to moment vectors in the interior of the moment space, where both principal representations exist and are unique; a vector on the boundary has exactly one representing measure, so the interval collapses to a point. $S=\kappa/(1-a-\kappa\Phi)$ is increasing in $\Phi$ on $\Phi<(1-a)/\kappa$, which gives the interval of summed responses, unbounded when $\Phi_+\ge(1-a)/\kappa$ because by convexity every $\Phi\in[\Phi_-,\Phi_+]$ is attained. (ii) If $\mu,\nu$ share moments up to order $n-1$, then $\int f\,d(\mu-\nu)=\int(f-p)\,d(\mu-\nu)$ for every polynomial $p$ of degree $n-1$, so $|\int f\,d(\mu-\nu)|\le2k_0\inf_p\|f-p\|_\infty=2k_0E_{n-1}$. With $r=\bar r(1+x)/2$, $f=(2/\bar r)/(x_0-x)$ on $[-1,1]$ with $x_0=2/\bar r-1>1$, and Chebyshev's formula gives $E_m(1/(x_0-x))=\varrho^m/(x_0^2-1)$ \citep{achieser1956}; we checked it against a linear-programming minimax fit for $m\le8$. Expanding $x_0=1+2(1-\bar r)/\bar r$ gives the asymptotics. (iii) Let the moment vector $c$ be interior to the moment cone of $[0,1]$. Interior points have an open neighbourhood in the cone, and $m(r_0)=(1,r_0,\ldots,r_0^{n-1})$ has norm at most $\sqrt n$, so for some $\epsilon>0$, $c-\epsilon m(r_0)$ is interior for every $r_0\in[0,1]$ and is represented by its lower principal representation, which has no atom at $1$. Adding $\epsilon\delta_{r_0}$ gives a measure on $[0,1)$ with moments $c$ and $\Phi\ge\epsilon/(1-r_0)\to\infty$ as $r_0\to1$; by convexity every larger $\Phi$ below $(1-a)/\kappa$ is attained by a stable economy, so $S$ is unbounded.

\paragraph{Theorem~\ref{thm:certcost}.} (i) The flow impulse enters the stacked state as $b=\kappa(1,\eta_1,\ldots,\eta_K)^\top$, and $h_t=e_x^\top J^tb\le\|J^tb\|_\infty\le\|J\|_\infty^t\|b\|_\infty$ because $J\ge0$. The price row of $J$ sums to $a+\kappa\Phi\le\bar s$ and cohort row $i$ to $1-\eta_i(1-a-\kappa\Phi)\le1-(1-\bar r)(1-\bar s)=\bar\lambda$, so $h_t\le\kappa\bar\lambda^t$ and $\sum_{t\ge T}h_t\le\kappa\bar\lambda^T/(1-\bar\lambda)$. $\sum_{t<T}\bar h_t$ is Gaussian with mean $\sum_{t<T}h_t$ and variance $T\sigma^2/N$, so $\sum_{t<T}\bar h_t+z_{1-\gamma}\sigma\sqrt{T/N}\ge\sum_{t<T}h_t$ with probability $1-\gamma$, and adding the tail bound gives validity. On the two-sided event of probability $1-2\gamma$, $\hat S^+-S\le2z_{1-\gamma}\sigma\sqrt{T/N}+\kappa\bar\lambda^T/(1-\bar\lambda)$; $T_\epsilon$ makes the second term at most $\epsilon/2$ and the stated $N$ makes the first at most $\epsilon/2$. (ii) Let $\theta_1$ raise the cohort weight of $\theta_0$ by $\Delta$ chosen so that $S_1-S_0=2\epsilon$; then $\|h_1-h_0\|_2=2C_0\epsilon(1+o(1))$. Validity at $\theta_1$ gives $P_1(\hat S\ge S_1)\ge1-\gamma$, hence $P_0(\hat S-S_0\ge2\epsilon)\ge1-\gamma-{\rm TV}(P_0,P_1)$. For $N$ independent Gaussian responses, ${\rm KL}=N\|h_1-h_0\|_2^2/(2\sigma^2)$, and Pinsker's inequality gives ${\rm TV}\le\sqrt N\,C_0\epsilon(1+o(1))/\sigma$. An overshoot probability below $1/2$ therefore needs ${\rm TV}>1/2-\gamma$, which is the stated bound. (iv) Write $M=\sum_iq_i(1-r_i)\delta_{r_i}$ for the merged measure, $k_m=\int r^m dM$ for the open-loop kernel and $\mathcal M_0$ for the measures on $[0,\bar r]$ of mass at most $B_0=(1-\bar r)(\bar s-a)/\kappa$; economies with merged measure in $\mathcal M_0$ lie in the class, and finite-atom measures are weak-$*$ dense in $\mathcal M_0$. \emph{Step 1 (closed loop to open loop).} The closed-loop transforms satisfy $H=\kappa/(1-az-\kappa zK)$, so $H_1-H_0=\kappa^{-1}zH_1H_0(K_1-K_0)$ and, by Young's inequality, $\|h_1-h_0\|_2\le\kappa^{-1}S_1S_0\|k_1-k_0\|_2\le C_1\|k_1-k_0\|_2$ on the class, for any number of recorded responses. \emph{Step 2 (two points).} If two economies have $S_1>S_0+\epsilon$, validity at the first and the overshoot bound at the second make $\{\hat S^+\le S_0+\epsilon\}$ have probability at least $1-\gamma$ under $P_0$ and at most $\gamma$ under $P_1$, so the total variation is at least $1-2\gamma$; for Gaussian responses this needs $\|h_1-h_0\|_2\ge2z_{1-\gamma}\sigma/\sqrt N$. Hence $\omega_S(2z_{1-\gamma}\sigma/(C_1\sqrt N))\le\epsilon$, where $\omega_S(\delta)$ is the largest difference in $S$ between two economies with $\|k_1-k_0\|_2\le\delta$. Since $\partial S/\partial\Phi=\kappa^2/(1-a-\kappa\Phi)^2\ge\kappa^2$, $\omega_S\ge\kappa^2\omega_\Phi$. \emph{Step 3 (duality).} For a depth $L$, $\omega_\Phi(\delta)$ over $\mathcal M_0$ with the first $L$ kernel values equals $\min_{c\in\mathbb R^L}\{\delta\|c\|_2+B_0\psi(g-p_c)\}$, with $g(r)=1/(1-r)$, $p_c(r)=\sum_{t<L}c_tr^t$ and $\psi(h)=(\sup h)_++(-\inf h)_+$ on $[0,\bar r]$: the inequality $\le$ is direct, and equality follows from Sion's minimax theorem because $\mathcal M_0^2$ is weak-$*$ compact \citep{donoho1994}; the modulus for all kernel values is the limit as $L\to\infty$, by compactness. \emph{Step 4 (a pole forces large coefficients).} If $\|c\|_2\le R$ and $\eta=\sup_{[0,\bar r]}|g-p_c|$, then $\log(1/\eta)\le K_0\max\{R^2,(1-\bar r)^{-1}\}$ with $K_0=8\log6\sqrt{1-\bar r}\log\rho_*$, $\rho_*=(1+\sqrt{1-\bar r})^2/\bar r$. Indeed $e=g-p_c$ is analytic on the disc of radius $\rho=1-s/2$, $s=\min\{1/(4R^2),(1-\bar r)/4\}$, with $|e(1-s)|\ge1/(2s)$ and $|e|\le3/s$ on $|z|=\rho$ because $|p_c(z)|\le R(1-|z|)^{-1/2}$; the two-constants theorem on the disc slit along $[0,\bar r]$, with the harmonic measure of the slit bounded below by an annulus comparison through the Joukowski map, gives the claim. \emph{Step 5.} For every $c$, either $\eta\ge\delta\sqrt{\log(1/\delta)}$, or $\log(1/\eta)>\tfrac12\log(1/\delta)$ and Step 4 forces $\|c\|_2\ge\sqrt{\log(1/\delta)/(2K_0)}$ once $\log(1/\delta)\ge2K_0/(1-\bar r)$; so $\omega_\Phi(\delta)\ge\min\{B_0,(2K_0)^{-1/2}\}\delta\sqrt{\log(1/\delta)}$ for small $\delta$, and Step 2 gives $N\ge c\,\sigma^2\epsilon^{-2}\log(1/\epsilon)$. (v) The identification rates are the companion paper's Theorems on the $K$-cohort lower bound, the matching upper bound and general memory.

\paragraph{Lemma~\ref{lem:sound}.} On the event that $\mathcal C(D)$ contains the true economy, the supremum over $\mathcal C(D)$ dominates the true worst-case loss for every action, including one selected from $D$; the event has probability at least $1-\gamma$.

\paragraph{Theorem~\ref{thm:robustguard}.} (i) For $c_I\ge0$, $\alpha\mapsto S(\alpha)=\kappa/(1-a-\kappa\alpha\Phi)$ is increasing, and so is the worst-case loss $|I|\bar uS(\alpha)$, which is also increasing in $\Phi$; the robust constraint is therefore the constraint at $\Phi^+$, and $|I|\bar u\kappa/(1-a-\kappa\alpha\Phi^+)\le\mathcal B$ is equivalent to $\alpha\le c_I/(\kappa\Phi^+)$. (ii) is Lemma~\ref{lem:sound}. (iii) At $\Phi=\Phi^+$ the constraint binds at $\alpha^*$ when $\alpha^*<1$. (iv) At $\alpha=0$, $S=\kappa/(1-a)$. (v) Upper: $\alpha^*$ is Lipschitz in $\Phi^+$ with constant $c_I/(\kappa\Phi_{\min}^2)$ and $\Phi^+$ is Lipschitz in $\hat S^+$ with constant $1/S_{\min}^2$, so Theorem~\ref{thm:certcost}(i) with $T\asymp\log N$ gives the rate. Lower: take two single-cohort economies in the class whose weights differ by $\Delta=c\,\sigma/\sqrt N$; their oracle capacities differ by a constant times $\Delta$, and ${\rm TV}\le\sqrt N\|\partial h/\partial q\|_2\Delta/(2\sigma)$, which is below $(1-\beta)/2$ for small $c$; Proposition~\ref{prop:twoworld} gives an expected loss of at least $(1-\beta)/2$ times the capacity difference.

\paragraph{Theorem~\ref{thm:disclosure}.} (i) If $O'$ reveals everything $O$ reveals, every economy compatible with $O'$ data is compatible with $O$ data, so suprema over $\Theta_{O'}$ are at most those over $\Theta_O$, and Theorem~\ref{thm:robustguard}(i) is monotone in $\Phi^+$. (ii) $\alpha^*$ is strictly decreasing in $\Phi^+$ when $c_I>0$ and $\alpha^*<1$, and $\Phi^+_{O'}<\Phi^+_O$ exactly when no economy with $\Phi=\Phi^+_O$ remains compatible (the supremum is attained, by Theorem~\ref{thm:window}(i)). (iii) follows from Theorem~\ref{thm:window}(iii). (iv) If $O$ identifies $\Phi$, $\Phi^+_O=\Phi$ and no contract can raise $\alpha^*$. (v) The delta method applied to $(w,r)\mapsto w/(1-r)$, whose gradient is $(1-r)^{-1}(1,\Phi_i)$, gives the standard error. If each certificate is valid at level $1-\gamma/2$, their minimum is valid at level $1-\gamma$ by the union bound and is no larger than either, so the admitted level is at least the larger of the two admitted levels.

\paragraph{Lemma~\ref{lem:certnum}.} For a compatible merged measure $M$ on $[0,\bar r]$, $\int p\,dM=\sum_jc_jk_j$ for any polynomial of degree at most $n-1$, so $p_-\le f\le p_+$ gives the two bounds by integration. Let $H$ interpolate $f$ at the atoms $x_i$ of $\mu_+$ with multiplicity $m_i=2$ at interior atoms and $1$ at $0$ and $\bar r$; principal representations have index $\sum_im_i=n$, so $H$ has degree at most $n-1$, and the Hermite remainder is $f(r)-H(r)=f^{(n)}(\xi)\,\omega(r)/n!$ with $\omega(r)=\prod_i(r-x_i)^{m_i}$. Since $f^{(n)}>0$ and, on $[0,\bar r]$, the factor $r-\bar r\le0$ appears once while every other factor is a square or $r\ge0$, $\omega\le0$ and $H\ge f$. For $\mu_-$ there is no atom at $\bar r$, so $\omega\ge0$ and $H\le f$. As $H=f$ on the support of $\mu_\pm$, $\sum_jc_jk_j=\int H\,d\mu_\pm=\int f\,d\mu_\pm=\Phi_\pm$. The certified computation proves $p_++s\ge f$ (or $p_--s\le f$) for the rounded polynomial and slack $s$ in interval arithmetic, in which every operation encloses its exact result \citep{rump2010}, so the first statement applies to the rounded polynomial with $c_0$ increased (decreased) by $s$.

\paragraph{Proposition~\ref{prop:hotelling}.} (i) The $X_e$ are i.i.d. normal with mean $S_T=\sum_{t<T}h_t$, whatever $\Sigma$, so $\sqrt N(\bar X-S_T)/s_X$ has Student's $t_{N-1}$ distribution and $P(S_T\le\bar X+t_{N-1,1-\gamma}s_X/\sqrt N)=1-\gamma$; the tail bound of Theorem~\ref{thm:certcost}(i) holds deterministically on the class. (ii) $N(\bar h-m)^\top\hat\Sigma^{-1}(\bar h-m)$ has the $\frac{T(N-1)}{N-T}F_{T,N-T}$ distribution \citep{hotelling1931}; on the event that it is at most $T^2_\gamma$, Cauchy--Schwarz in the $\hat\Sigma$ inner product gives $c^\top(m-\bar h)\le\sqrt{c^\top\hat\Sigma c}\,\sqrt{(\bar h-m)^\top\hat\Sigma^{-1}(\bar h-m)}\le\sqrt{T^2_\gamma c^\top\hat\Sigma c/N}$ for every $c$ at once. (iii) If the taps of an experiment have common variance $\sigma^2$ and correlation matrix $R$, the margin $z_{1-\gamma}\sigma\sqrt{T/N}$ has coverage $\Phi\bigl(z_{1-\gamma}\sqrt{T/\mathbf 1^\top R\mathbf 1}\bigr)$, below $1-\gamma$ whenever $\mathbf 1^\top R\mathbf 1>T$; with $T=12$ and autocorrelation $0.8$ it is $0.751$, matching the simulations.

\paragraph{Proposition~\ref{prop:venue}.} (i) $D$ is a quadratic with leading coefficient $\delta\rho\ge0$, and a level is safe if and only if $D(\alpha)\ge\kappa_0|I|\bar u/\mathcal B$ (which implies $D>0$, hence stability and $S=\kappa_0/D$). The set where a convex function is below a constant is an interval, so the safe set is $[0,1]$ minus an open interval; the example in the text has two components. (ii) For $\alpha\ge0$ and $a_0<1$, $\partial D/\partial\delta=\alpha(1-a_0+\rho\alpha)\ge0$, $\partial D/\partial\rho=\alpha(1+\delta\alpha)\ge0$, $\partial D/\partial a_0=-(1+\delta\alpha)\le0$ and $\partial D/\partial\Phi=-\kappa_0\alpha\le0$, while the condition $D-\kappa_0|I|\bar u/\mathcal B\ge0$ decreases in $\kappa_0$. The minimum of the left side over the box is therefore at the stated corner for every $\alpha$. (iii) follows from (i) and (ii).

\paragraph{Theorem~\ref{thm:venuebb}.} (i) In interval arithmetic with outward rounding every operation returns an interval containing the exact result for all real arguments in its operands \citep{rump2010}; by induction the computed interval vectors contain $J(\alpha,\theta)^ts_0(\alpha,\theta)$ and the interval matrices contain $J^j$ for every $(\alpha,\theta)$ in the box, so the computed sum contains $\sum_{t<T}|h_t|$. For $t\ge T$, writing $t=T+kT+j$ with $k\ge0$ and $0\le j<T$, $|h_t|\le\|J^{kT+j}J^Ts_0\|_\infty\le\|J^T\|_\infty^k\|J^j\|_\infty\|J^Ts_0\|_\infty$, and summing gives the tail bound. A level interval classified safe has an upper enclosure within budget over every part of a partition of $\Theta$, hence for every $\theta$; one classified unsafe has $S(\alpha,\theta_j)\ge\sum_{t<T}|h_t(\alpha,\theta_j)|>$ budget for some $\theta_j\in\Theta$ and every level in it. (ii) Each $S(\cdot,\theta)$ is continuous where finite and tends to infinity at the stability boundary, so $\alpha\mapsto\sup_\Theta S(\alpha,\theta)$ is lower semicontinuous and its sublevel set, the robust safe set, is closed; the brackets follow from (i). (iii) is immediate. (iv) Under $\sup_\Theta\rho(J)<1$ there is $T$ with $\|J^T\|_\infty<1$ uniformly on $\Theta$; interval extensions of rational functions converge linearly in the box width \citep{hansen2004}, so at a level where the robust risk differs from the budget by $g>0$ the enclosures over small enough boxes, with $T$ large enough that the truncation is below $g/2$, decide.

\paragraph{Proposition~\ref{prop:probe}.} With $D_n=\sum_{e\le n}(X_e-S_T)$ and $V_n=n\sigma_T^2$, $\exp(\lambda D_n-\lambda^2V_n/2)$ is a nonnegative supermartingale for every $\lambda$ under sub-Gaussian increments; mixing over $\lambda\sim N(0,1/\varrho)$ gives the nonnegative supermartingale $M_n=\sqrt{\varrho/(V_n+\varrho)}\exp\bigl(D_n^2/(2(V_n+\varrho))\bigr)$ with $M_0=1$. Ville's inequality \citep{ville1939} gives $P(\exists n:M_n\ge1/\gamma)\le\gamma$, and $M_n\ge1/\gamma$ is exactly $|D_n|\ge b_n$. On the complementary event every interval covers $S_T$ at every $n$, so a decision made at any stopping time is correct; the tail bound is deterministic. Since $b_n/n=O\bigl(\sigma_T\sqrt{\log(n/\gamma)/n}\bigr)$, the stopping time is $O\bigl(\sigma_T^2\Delta^{-2}\log(\sigma_T/(\Delta\gamma))\bigr)$ when the gap is $\Delta$ \citep{howard2021}.

\section{Notation}\label{app:notation}
\begin{longtable}{p{0.2\textwidth}p{0.7\textwidth}}
\toprule
Symbol & Meaning\\
\midrule
\endhead
$x_t,\ d_t,\ u_t$ & price displacement, depth imbalance $(B-A)/(B+A)$, signed flow\\
$a,\ \kappa$ & price persistence and direct impact\\
$\varphi,\ \varepsilon,\ r$ & price sensitivity to depth, depletion per unit of flow, refill rate\\
$\eta,\ q$ & momentum memory rate and aggressiveness\\
$C_2,\ C_3$ & amplification statistics $\varphi\varepsilon-(1-a)\kappa$ and $C_2+\kappa^2\eta q$\\
$h_k,\ g_k$ & impulse response of the book; price impulse coefficient of the controlled system\\
$J(\alpha)$ & closed-loop operator at admission level $\alpha$\\
$M_C,\ G_{\rm cum}$ & peak and summed price-row gains\\
$\mathcal K(J)$ & Kreiss constant\\
$\bar u,\ \bar u_{\rm step},\ \bar u_{\rm cum}$ & shock envelopes\\
$B_{\rm rem},\ R_{\rm safe},\ \mu$ & remaining margin buffer, safety reserve, maintenance-margin rate\\
$\epsilon,\ S$ & operator confidence radius and $\sum_k\|J^k\|$\\
$M_{\rm std},\ A_{xx}$ & unit-free and price-to-price spot amplification statistics\\
\bottomrule
\end{longtable}

\section{Additional material}\label{app:ia}

\subsection{Specification of the pre-registered tests}
\paragraph{Futures (Table~\ref{tab:d8}).} Question: does $\hat C_3$, estimated from the previous six hours of futures depth, flow and price, predict the next hour's MAE beyond controls? Unit: hourly blocks 06--23 UTC, features from $[t-360,t)$ minutes, outcomes from $(t,t+60]$. Controls: $rv_{60}$, $rv_{360}$, $|r_{60}|$, flow imbalance, log volume, log depth, imbalance, open-interest change. Model: OLS; tail is the outcome above the training 90th percentile; score is the fitted value. Inference: Clark--West with Newey--West (6 lags); day-block bootstrap with 2,000 draws. Secondary: peak ratio as candidate; open-interest drop and depth collapse as outcomes; Holm. Train $\le$ 2024-12-31; test $\ge$ 2025-01-01. Hash \texttt{d95212187ae42052591a5ca5acf73317} \texttt{0f217f430a5a6ca812da6e48a1e4af85}; one unlock.

\paragraph{Spot (Section~\ref{sec:spot}).} Question: does pre-window $A_{xx}$ predict the next ten minutes' MAE beyond volatility and flow controls? Unit: non-overlapping 30-minute blocks; features from the previous ten minutes. Controls: $rv_{10m}$, $rv_{60m}$, $|r_{10m}|$, flow imbalance, log volume, spectral radius, hour of day. Train 2021--2023; test 2024. Hash \texttt{b1fd7f8eef15df1ba295eb6d32b53be6} \texttt{c23ea2d23c5216de912555e62716e5ae}; three unlocks (see text).

\subsection{Adaptive safe probing}\label{app:probe}
The certificate of Theorem~\ref{thm:certcost} fixes the number of reset experiments in advance, so it must be sized for the hardest economy the supervisor wants to decide. Experiments can instead be run one at a time and stopped as soon as the data decide. Let each experiment return $X_e=\sum_{t<T}h^{(e)}_t$, independent with mean $S_T$ and sub-Gaussian proxy $\sigma_T^2$ ($\sigma_T^2=T\sigma^2$ under (O1)), and let $b_n=\sqrt{(n\sigma_T^2+\varrho)\bigl(\log((n\sigma_T^2+\varrho)/\varrho)+2\log(1/\gamma)\bigr)}$ for a fixed $\varrho>0$.

\begin{proposition}[Anytime-valid probing]\label{prop:probe}
With probability at least $1-\gamma$, $|\bar X_n-S_T|<b_n/n$ for every $n$ simultaneously. Hence the rule that probes until $\bar X_n+b_n/n+\kappa\bar\lambda^T/(1-\bar\lambda)\le s_{\mathcal B}$ (admit), or $\bar X_n-b_n/n>s_{\mathcal B}$ (refuse), or a maximum number of probes is reached (abstain), makes a wrong decision with probability at most $\gamma$ whatever its stopping time, where $s_{\mathcal B}=\mathcal B/(|I|\bar u)$. When $S_T$ is at distance $\Delta$ from $s_{\mathcal B}$ it stops after $O\bigl(\sigma_T^2\Delta^{-2}\log(\sigma_T/(\Delta\gamma))\bigr)$ probes. Each probe is itself admissible if its position $v_p$ satisfies $|v_p|\bar u_p\kappa/(1-\bar\lambda)\le\mathcal B_p$, the worst case over the class.
\end{proposition}

The value is in adapting to the economy. With $\sigma_T=0.17$ and a threshold of $2.0$, the sequential rule decides economies at distance $0.5$ from the threshold after 3.7 probes on average, at distance $0.2$ after 11, at $0.1$ after 30 and at $0.05$ after 107, with no wrong decision in 7{,}000 runs, while a fixed design that must decide at distance $0.05$ uses 130 probes for every economy (test T14). Probes have a monetary cost and move prices, so a supervisor who can probe sequentially spends where the decision is hard. This is an extension rather than a result the paper relies on; optimal probe design under impact costs is left open.

\section{Reproducibility: source of every number}\label{app:repro}

The replication package contains numbered Colab notebooks and stand-alone scripts; its README lists them in run order. ``Part $N$'' below and in table notes refers to notebook $N$ (Part 1: theorem checks and spot tests; Part 2: BTCUSDT futures depth; Parts 4--18: the cross-section, replications, fresh-day tests, second venues, the CFTC test, the certified guard, post-hoc checks, US equities, the information-safety tests on untouched days and symbols, the new-venue test and the last post-hoc checks). Result identifiers such as \texttt{X6} or \texttt{U6} name the JSON file that the notebook writes. ``Kind'' distinguishes closed-form or simulation results (recomputed), results from public market data (real data), and audits.

\begin{longtable}{p{0.37\textwidth}p{0.37\textwidth}p{0.14\textwidth}}
\toprule
Result & Source & Kind\\
\midrule
\endhead
Spectral blindness gains & \texttt{T5\_spectral\_blindness} & recomputed\\
Table~\ref{tab:kreiss}; burst operator gain, $q$, certificate & \texttt{kreiss\_table.py} $\to$ \texttt{R\_P2\_kreiss.json} & recomputed\\
Table~\ref{tab:calib}, Figure~\ref{fig:calib} & \texttt{calib\_guard.py} $\to$ \texttt{R\_P2\_calibrated\_guard.json} & recomputed\\
Figure~\ref{fig:pseudo} & \texttt{make\_figs2.py} & recomputed\\
Theorem~\ref{thm:peak}\ref{thm:t6}, Theorem~\ref{thm:peak}\ref{prop:shape}--\ref{prop:t6m} draws & \texttt{T6\_amplification\_theorem} & recomputed\\
Guard maximality; peak guard unsafe; monotonicity & \texttt{T7\_guard\_maximality}, \texttt{P2.11-T1} & recomputed\\
Robust gain, 30,000 draws & \texttt{T8\_robust\_gain} & recomputed\\
Operator ball coverage; certificate share & \texttt{P2.4-T1\_coverage} & recomputed\\
Summed gain vs.\ weights & \texttt{T9\_bridge} & recomputed\\
Depth validation (Table~\ref{tab:d2}) & \texttt{D2\_depth\_validation} & real data\\
Table~\ref{tab:d3} & \texttt{D3\_T6\_structural\_train} & real data\\
Figure~\ref{fig:med}, Table~\ref{tab:d4} & \texttt{D4\_depth\_mediation\_train} & real data\\
No-look-ahead audit & \texttt{D5\_no\_lookahead} & audit\\
Table~\ref{tab:d8}; hash & \texttt{D8\_depth\_test\_G2}, \texttt{P2\_depth.json} & real data\\
Flag comparison & \texttt{D8b\_flag\_comparison} & real data\\
Table~\ref{tab:d9}, Figure~\ref{fig:events} & \texttt{D9\_event\_studies} & real data\\
Unit dependence $8.73\to86{,}731$ & \texttt{P2.2-T1\_unit\_invariance} & real data\\
Table~\ref{tab:spotnull}; sensitivity & \texttt{P2.2-T2}, \texttt{T3}, \texttt{T4}, \texttt{T5} & real data\\
Table~\ref{tab:spotlp} & \texttt{T6-E\_overshoot\_real} & real data\\
Spot predictive test & \texttt{P2.5\_predictive\_G2}, \texttt{P2.5-T2}, \texttt{T3} & real data\\
Spot flag comparison & \texttt{P2.7-T2\_vol\_ablation} & real data\\
Theorem~\ref{thm:sum}; Remark~\ref{rem:counter}; Theorem~\ref{thm:peak}\ref{thm:nomom} check & \texttt{equilibrium.py}, \texttt{t6m\_search*.py}, \texttt{thmA\_check.py} & recomputed\\
Proposition~\ref{prop:equil} numbers & \texttt{equilibrium.py} $\to$ \texttt{R\_P2\_equilibrium.json} & recomputed\\
Proposition~\ref{prop:invariance} check; multi-asset example; Proposition~\ref{prop:lucas} check & \texttt{recompute\_oct20.py} & recomputed\\
Tables~\ref{tab:x3}, \ref{tab:x4}, \ref{tab:x8}; tick events & Part 4: \texttt{X3}, \texttt{X4}, \texttt{X8}, \texttt{X7} & real data\\
Table~\ref{tab:x6}; flags & Part 4: \texttt{X6\_xsec\_test\_G2}, \texttt{X6b}; \texttt{P4\_xsec.json} (5910b6b4\ldots) & real data\\
Table~\ref{tab:x6r} & Part 4: \texttt{X6r\_robustness\_posthoc} & real data (post hoc)\\
Table~\ref{tab:y6} & Part 5: \texttt{Y6\_replication\_G2}; \texttt{P5\_replication.json} (6aa8e2c0\ldots) & real data\\
Table~\ref{tab:y8} & Part 5: \texttt{Y8\_kappa\_control\_posthoc} & real data (post hoc)\\
Tables~\ref{tab:z6}, \ref{tab:backtest}; equilibrium check & Part 6: \texttt{Z6\_*}, \texttt{Z8}; \texttt{P6\_fresh\_days.json} (65acda36\ldots) & real data\\
Table~\ref{tab:certguard}, Figure~\ref{fig:oct10}, economic size (Section~\ref{sec:threats}) & Part 12: \texttt{Q6\_certified\_guard}, \texttt{Q7\_oct10\_event}, \texttt{E1\_economic\_size}; \texttt{P12\_certified\_guard.json} (b9a56af4\ldots) & real data\\
Liquidation probabilities and dollar losses (Section~\ref{sec:threats}) & \texttt{recompute\_dollar.py} (from \texttt{E1\_economic\_size}) $\to$ \texttt{R\_P2\_dollar\_translation.json} & recomputed\\
Table~\ref{tab:bybitob} & Part 10: \texttt{V3A\_orderbook\_download}, \texttt{V4\_build}, \texttt{V6\_bybit\_ob\_test}; \texttt{P10\_bybit\_ob.json} (ec2145fa\ldots) & real data\\
Table~\ref{tab:fixedgrid}; throttle, Bybit-bar, instrument, tercile and VPIN checks & Part 13: \texttt{R5}, \texttt{R1}, \texttt{R3}, \texttt{R4}, \texttt{R2}, \texttt{R8} & real data (post hoc)\\
Example~\ref{ex:witness}; Theorems~\ref{thm:window}, \ref{thm:certcost}, \ref{thm:robustguard}, \ref{thm:disclosure} numbers; Tables~\ref{tab:certvsid}, \ref{tab:discsim}, \ref{tab:simtests} & \texttt{run\_all.py} in \texttt{research\_extension/} \texttt{information\_safety} $\to$ \texttt{IS\_results.json}; unit tests & recomputed\\
Section~\ref{sec:p15} & Part 15: \texttt{W4\_build}, \texttt{W5\_blocks}, \texttt{W6\_info\_safety\_test}; \texttt{P15\_info\_safety.json} (\texttt{69849c96}) & real data\\
Section~\ref{sec:p16} & Part 16: \texttt{X2\_new\_universe}, \texttt{X4\_build}, \texttt{X5\_blocks}, \texttt{X6\_untouched\_symbols\_test}; \texttt{P16\_untouched\_symbols.json} (\texttt{0523e476}) & real data\\
Section~\ref{sec:p17} & Part 17: \texttt{C2\_universe}, \texttt{C4\_build}, \texttt{C5\_blocks}, \texttt{C6\_new\_venue\_test}; \texttt{P17\_new\_venue.json} (\texttt{3233c52b}) & real data\\
Section~\ref{sec:p18} & Part 18: \texttt{R3\_placebo\_binance}, \texttt{R4\_guard\_baselines}; \texttt{P18\_posthoc\_plan.json} & real data (post hoc)\\
Section~\ref{sec:p19} (Table~\ref{tab:p19}) & Part 19: \texttt{recompute\_\allowbreak p2\_\allowbreak parts19\_\allowbreak 20\_\allowbreak 21.py} $\to$ \texttt{R\_P2\_\allowbreak parts19\_\allowbreak 20\_\allowbreak 21.json} & real data\\
Section~\ref{sec:p20} (Table~\ref{tab:p20}) & Part 20: \texttt{recompute\_\allowbreak p2\_\allowbreak parts19\_\allowbreak 20\_\allowbreak 21.py} $\to$ \texttt{R\_P2\_\allowbreak parts19\_\allowbreak 20\_\allowbreak 21.json} & real data\\
Section~\ref{sec:p21} (Table~\ref{tab:p21}) & Part 21: \texttt{recompute\_\allowbreak p2\_\allowbreak parts19\_\allowbreak 20\_\allowbreak 21.py} $\to$ \texttt{R\_P2\_\allowbreak parts19\_\allowbreak 20\_\allowbreak 21.json} (from \texttt{W6}, \texttt{X6}) & translation of real data\\
Sections~\ref{sec:p22}--\ref{sec:p23} (Tables~\ref{tab:p22}, \ref{tab:p23}) & Parts 22--23: \texttt{recompute\_\allowbreak prereg\_\allowbreak oct21\_\allowbreak 22nd.py} $\to$ \texttt{R\_P2\_\allowbreak p22\_\allowbreak p23.json}; \texttt{P22\_\allowbreak liquidation\_\allowbreak mechanism.json} (\texttt{2b3794e3}), \texttt{P23\_\allowbreak guard\_\allowbreak capacity.json} (\texttt{d4169377}) & real data\\
Sections~\ref{sec:p24}--\ref{sec:p25} (Tables~\ref{tab:p24}, \ref{tab:p25}) & Parts 24--25: \texttt{recompute\_\allowbreak prereg\_\allowbreak oct21\_\allowbreak followup.py} $\to$ \texttt{R\_P2\_\allowbreak p24\_\allowbreak p25.json}; \texttt{P24\_\allowbreak guard\_\allowbreak throttle\_\allowbreak and\_\allowbreak capacity.json} (\texttt{698f9119}), \texttt{P25\_\allowbreak persistence\_\allowbreak mechanism.json} (\texttt{8ebd67f3}) & real data\\
Inventory and corrected hidden-state checks & \texttt{recompute\_\allowbreak p2\_\allowbreak step2\_\allowbreak theory.py} $\to$ \texttt{R\_P2\_\allowbreak step2\_\allowbreak theory\_repaired.json} & recomputed\\
Balance-sheet Monte Carlo (Section~\ref{sec:guard}) & \texttt{bridge/\allowbreak microfoundation\_\allowbreak and\_\allowbreak baselines.py} $\to$ \texttt{bridge/\allowbreak results/\allowbreak microfoundation\_\allowbreak and\_\allowbreak baselines.json} & recomputed\\
Lemma~\ref{lem:certnum}; Propositions~\ref{prop:hotelling}, \ref{prop:venue}, \ref{prop:probe}; Theorem~\ref{thm:venuebb} & \texttt{run\_all.py}: \texttt{T10}, \texttt{T11}, \texttt{T13}, \texttt{T14} & recomputed (certified for T10, T13)\\
Table~\ref{tab:us} & Part 14: \texttt{A1\_us\_universe}, \texttt{A3\_us\_build}, \texttt{A4\_us\_blocks}, \texttt{A5\_us\_equity\_test}; \texttt{P14\_us\_equities.json} (f5800a69\ldots) & real data\\
Theorem~\ref{thm:peak}\ref{prop:shape} peak bracket; Remark~\ref{rem:counter}; Theorem~\ref{thm:sum}, \ref{prop:equil} checks & \texttt{recompute\_audit.py} $\to$ \texttt{R\_audit\_oct2026.json} & recomputed\\
Theorem~\ref{thm:peak} residues and confluent peak lag; Remark~\ref{rem:cheb}; Proposition~\ref{prop:lucas}(ii) & \texttt{hardening/p2\_sec4\_checks.py}, \texttt{p2\_cheb\_tradeoff.py}; unit tests \texttt{test\_hardening\_theory.py} & recomputed\\
Theorem~\ref{thm:disclosure}(v), intersection row of Table~\ref{tab:discsim} & information-safety package, test T15 & recomputed\\
Section~\ref{sec:bybit} & Part 7: \texttt{W6\_bybit\_test}; \texttt{P7\_bybit.json} & real data\\
Table~\ref{tab:z9} & Part 9: \texttt{U6\_volbench\_test}; \texttt{P9\_volbench.json} (71983f1a\ldots) & real data\\
Figure~\ref{fig:forest} & \texttt{make\_fig\_forest.py} & real data\\
\bottomrule
\end{longtable}





% BEGIN ECONOMIC REVIEW 20261029
\section*{Development revision and validation status: 29 October 2026}
This is a working draft for author and specialist review. The revision corrects
hidden-state drift accounting, distinguishes cumulative displacement from peak
mark-to-market loss, and restricts claims to their stated model and observation
contracts. It does not declare every theorem independently verified or every
research objective complete.

\paragraph{Returned execution record.} The author returned a completed full-mode
notebook archive using source commit \texttt{19805d7e}, seed 20261029,
Python 3.13.18 and NumPy 2.1.3. All 10 exported-file hashes and the four analytical
source modules were checked. The run completed E1, E2, E3, E4 and E6; its log
reports 18 targeted tests passing. This is an execution reproduction of the same
seeded code, not independent scientific validation. No real-market CSV was
provided, and the GPU, known-source-feedback and prospective studies were not run.

\paragraph{What the checks establish.} E1 cross-checks the corrected drift identity
against a resolvent and exercises a sufficient invariant-box guard. E2 demonstrates
filter-induced persistence under martingale nulls. E3 compares known alternatives
with equal aggregate risk; it is an illustrative witness, not a uniform rate
theorem. E4 matches baseline capacity on training episodes and reports test
capacity differences; the guard and volatility control both have zero observed
budget breaches in the three bounded-shock synthetic regimes. This does not
establish universal guard superiority or economic welfare.

\paragraph{Coverage and the observation contract.} E6 uses 1,000 replications,
32 reset responses and 512 taps. The nominal 95\% interval assumes independent
Gaussian tap errors and a valid memory envelope. The intervals in the table
measure Monte Carlo uncertainty about coverage, not financial risk.
\begin{center}\small
\begin{tabular}{lrrc}
\toprule
Setting & Coverage & MC 95\% interval & Contract holds\\
\midrule
Independent Gaussian & 94.6\% & [93.0, 95.8] & yes\\
Correlated tap errors & 48.3\% & [45.2, 51.4] & no\\
Heavy-tailed readings & 95.1\% & [93.6, 96.3] & no\\
Biased readings & 25.6\% & [23.0, 28.4] & no\\
Unseen slow memory & 0.0\% & [0.0, 0.4] & no\\
\bottomrule
\end{tabular}
\end{center}
Failures under correlated errors, biased readings and a false memory floor expose
the contract's limits. The heavy-tail outcome in this configuration is not a
general robustness result. All outcomes, including negative ones, are retained.

\paragraph{Remaining obligations.} Independent review of the load-bearing
identification and certification proofs, continuous numerical extrema, economic
primitives, the public-observation bridge and policy counterfactuals remains
necessary. The current two-paper name-map audit finds 54 of 57 formal environments
with existing declaration mappings, three unmapped statements and seven revised
statements requiring semantic review. A label mapping is not a new Lean proof.
The experiment design and full results are retained under
\texttt{studies/economic\_review\_20261029/}.

% END ECONOMIC REVIEW 20261029

% BEGIN INFORMATION DECISION REVISION 20261029
\section*{Information-supported portfolio choice: economic revision and returned evidence}

\subsection*{What this revision adds}
The information question is distinct from choosing a preferred amount of risk.
For a specified financial decision, the relevant issue is which observations
support an exposure constraint, what remains unidentified, and whether obtaining
additional information changes the attainable decision after costs. This does
not introduce portfolio choice or control without full identification as new
ideas. The general identification-versus-decision-value result still requires
independent proof and a comparison with the closest existing results.

\subsection*{Time conventions and the institution}
The new application is a price-taking, unlevered own-account investor, not a
venue-wide supervisor. It begins with cash, buys once, holds a fixed position
for 60 one-minute model periods, and closes fully. Round-trip linear and
quadratic charges enter cash and wealth. The account cannot change the market's
feedback, and the experiment does not measure customer liquidity or welfare.
This is separate from the paper's dealer/venue admission contracts.

\begin{center}\small
\begin{tabular}{p{0.26\textwidth}p{0.65\textwidth}}
\toprule
Quantity & Interpretation\\
\midrule
Sampling interval & One minute in the new synthetic account; other historical studies retain their stated sampling clocks.\\
Historical lookback & Often six hours in the existing forecasting studies; not a holding period.\\
Forecast horizon & Often the next hour's adverse excursion; not an executed round-trip trade.\\
New holding period & Exactly 60 model minutes, with explicit entry and exit charges.\\
Other theory horizons & The existing finite-horizon theorems retain their declared $H$; no universal holding period is imposed.\\
\bottomrule
\end{tabular}
\end{center}

\subsection*{Information and decision design}
The new comparison uses controlled independent reset observations of 12 impulse
taps, a known plant and reading variance, and a declared finite candidate class.
It does not treat public market trading as a reset experiment. Common training
and validation paths supply 768 and 384 additional reset episodes per repetition;
512 further paths evaluate the selected actions. These are in addition to the
reported tap-reading budget, not evidence of learning from 32 observations in
total. Verified disclosure restricts
the minimum-memory rate to an audited band. The oracle knows the true response,
not future shocks. All certificate arms use the same training objective.

The account selects a position by empirical expected logarithmic terminal
wealth, subject to its information-supported exposure restriction. This is a
finite-horizon application of logarithmic utility, not a newly proved Kelly
growth-optimality theorem. Equal response bounds alone do not imply equal
expected returns or identical utility-optimal portfolios.

The new account has initial wealth of 1,000 dollars and an initial price of
100 dollars; its position grid runs from zero to 9.9 units. Its round-trip
dollar charge is $c(I)=0.002I+0.0002I^2$. For bounded external flow, the
constraint adds this charge to position times the largest admissible absolute
price response over the holding horizon. The resulting bound concerns loss
from initial wealth at each date. It is not peak-to-trough drawdown, protection
against arbitrary unbounded shocks, or a rolling deployment guarantee.

\subsection*{Returned execution and proof scope}
The returned full-mode archive uses seed 20261029. All 35 manifest entries and
eight core source files were checked. The run contains 20 independent
repetitions per regime, not 16,800 independent experiments: the 16,800 policy
rows reuse tapes across policies, budgets and observation counts. Twelve
targeted checks passed. Eight conditional decision-logic lemmas compiled with
Lean 4.19.0 and their axiom printouts report no axioms. This is a seeded
execution reproduction. It neither completes the older manuscript's Lean
coverage nor formalizes the Gaussian confidence argument, market dynamics,
general minimax separation, portfolio optimality or welfare.

The displayed slice uses 32 reset episodes and a one-dollar loss budget from
initial wealth of 1,000 dollars. It is an illustrative presentation of an
exploratory grid, not a primary selected before the results were seen. All
regimes, budgets, policies and negative outcomes remain in the result archive.
The inputs and complete outputs are retained with the verification record in
the new information-decision study directory.

\subsection*{Aggregate information and unresolved cohorts: returned simulation}
\begin{center}\small
\begin{tabular}{lrrr}
\toprule
Regime & Unique candidate (\%) & Certified units & Oracle units\\
\midrule
No feedback & 0.0 & 9.405 & 9.900\\
Moderate feedback & 5.0 & 5.000 & 5.000\\
Near boundary & 0.0 & 1.200 & 1.200\\
\bottomrule
\end{tabular}
\end{center}
Here ``unique'' means one remaining candidate in the declared finite class,
not identification of arbitrary real-world algorithms. In the moderate and
near-boundary configurations, the aggregate certificate reaches the oracle's
grid capacity in this slice even though the candidate decomposition is usually
unresolved. This is an illustrative decision comparison; it does not prove a
uniform rate separation for continuous classes.

\subsection*{Financial decisions: returned simulation}
{\footnotesize\setlength{\tabcolsep}{4pt}\renewcommand{\arraystretch}{0.9}
\begin{longtable}{p{0.36\textwidth}rrrr}
\toprule
Policy & Units & Log growth (bp) & ES95 (\$) & Breach (\%)\\
\midrule
\endfirsthead
\toprule
Policy & Units & Log growth (bp) & ES95 (\$) & Breach (\%)\\
\midrule
\endhead
\multicolumn{5}{l}{\textit{No feedback}}\\
Aggregate certificate & 9.405 & 0.578 & 0.394 & 0.000\\
Unpaid disclosure & 9.405 & 0.578 & 0.394 & 0.000\\
Paid disclosure & 9.405 & 0.478 & 0.404 & 0.000\\
Oracle response & 9.900 & 0.614 & 0.415 & 0.000\\
Volatility and persistence & 9.900 & 0.614 & 0.415 & 0.000\\
Point response & 9.900 & 0.614 & 0.415 & 0.000\\
Expected-log objective & 9.900 & 0.614 & 0.415 & 0.000\\
Fractional log, validation & 9.900 & 0.614 & 0.415 & 0.000\\
Empirical loss constraint & 9.900 & 0.614 & 0.415 & 0.000\\
Robust empirical log & 0.000 & 0.000 & 0.000 & 0.000\\
\midrule
\multicolumn{5}{l}{\textit{Moderate feedback}}\\
Aggregate certificate & 5.000 & 1.330 & 0.234 & 0.000\\
Unpaid disclosure & 5.000 & 1.330 & 0.234 & 0.000\\
Paid disclosure & 4.900 & 1.205 & 0.240 & 0.000\\
Oracle response & 5.000 & 1.330 & 0.234 & 0.000\\
Volatility and persistence & 8.550 & 2.214 & 0.404 & 0.000\\
Point response & 5.000 & 1.330 & 0.234 & 0.000\\
Expected-log objective & 9.900 & 2.537 & 0.469 & 0.000\\
Fractional log, validation & 9.900 & 2.537 & 0.469 & 0.000\\
Empirical loss constraint & 9.900 & 2.537 & 0.469 & 0.000\\
Robust empirical log & 5.195 & 1.374 & 0.244 & 0.000\\
\midrule
\multicolumn{5}{l}{\textit{Near boundary}}\\
Aggregate certificate & 1.200 & 1.386 & 0.084 & 0.000\\
Unpaid disclosure & 1.200 & 1.386 & 0.084 & 0.000\\
Paid disclosure & 1.200 & 1.286 & 0.094 & 0.000\\
Oracle response & 1.200 & 1.386 & 0.084 & 0.000\\
Volatility and persistence & 2.555 & 2.943 & 0.180 & 0.000\\
Point response & 1.200 & 1.386 & 0.084 & 0.000\\
Expected-log objective & 9.900 & 11.252 & 0.705 & 0.195\\
Fractional log, validation & 9.900 & 11.252 & 0.705 & 0.195\\
Empirical loss constraint & 9.900 & 11.252 & 0.705 & 0.195\\
Robust empirical log & 8.730 & 9.941 & 0.620 & 0.078\\
\midrule
\bottomrule
\end{longtable}}
Values average 20 repetitions. Log growth is $10^4$ times the mean logarithm
of terminal wealth divided by initial wealth; ES95 is the mean of each
repetition's worst five percent of initial-wealth losses. It is not a return
forecast or a peak-to-trough drawdown guarantee. Position sizes differ across
policies; the table reports a risk--growth tradeoff rather than equal-exposure
dominance. A zero observed breach is not proof of zero breach probability.

\subsection*{Disclosure cost and the negative result}
\begin{center}\small
\begin{tabular}{lrrr}
\toprule
Paid minus aggregate & Mean (bp) & Lower (bp) & Upper (bp)\\
\midrule
No feedback & -0.100 & -0.100 & -0.100\\
Moderate feedback & -0.126 & -0.126 & -0.125\\
Near boundary & -0.100 & -0.100 & -0.100\\
\bottomrule
\end{tabular}
\end{center}
The intervals are exploratory percentile bootstraps over repetition-level
paired means, not simultaneous intervals across the grid. In this slice the
unpaid disclosure does not expand capacity, while the assumed one-cent charge
reduces net log wealth. This run therefore does not establish a positive net
economic value of disclosure. The price is hypothetical, not a calibrated
reporting cost. Stronger disclosure need not be worth purchasing when aggregate
information already supports the relevant decision.

\subsection*{Assumption failures and calibration diagnostics}
The table reports finite-sample frequencies for all seven regimes, combining
the three reading budgets. Each band outcome is counted once per repetition
and reading budget (60 outcomes per regime); repeating the same outcome across
four financial loss budgets would not create further observations. The three
reading budgets within a repetition are not an independent replication.
\begin{center}\small
\begin{tabular}{lcrr}
\toprule
Regime & Contract valid & Tap-band hit (\%) & Empty set (\%)\\
\midrule
No feedback & Yes & 96.7 & 3.3\\
Moderate feedback & Yes & 93.3 & 5.0\\
Near boundary & Yes & 96.7 & 3.3\\
Hidden slow memory & No & 98.3 & 11.7\\
Correlated readings & No & 98.3 & 1.7\\
Biased readings & No & 0.0 & 100.0\\
Behaviour change & No & 95.0 & 5.0\\
\bottomrule
\end{tabular}
\end{center}
The nominal one-decision band coverage is 95 percent under the valid contract;
these small, dependent-budget frequencies are diagnostics, not a new coverage
theorem. A tap-band hit concerns the observed response, not membership in the
finite candidate class or stability of future behaviour. In particular, hidden
slow memory falls outside the declared class even when its first 12 taps fit
the bands. Under behaviour change the readings describe the old response.
Biased readings empty the set throughout this run and force the certificate
policy to abstain. Correlated readings use an unverified standard-error model.
None of these four regimes retains the stated guarantee; a high observed band
hit or zero observed financial breaches does not make an invalid contract valid.

\subsection*{Pending evidence: deliberately empty result cells}
The following numerical cells remain empty. The returned archive contained no
market CSV, dynamic-inventory experiment, institution-level intervention or
external replication. The general theorem and independent novelty review also
remain open; a table placeholder is not a theorem statement or a proof.

\begin{center}\small
\begin{tabular}{p{0.62\textwidth}cc}
\toprule
Pending study & Estimate & Uncertainty\\
\midrule
Untouched market-data test beyond order-flow imbalance and depth & & \\[1.8ex]
Full changing-inventory, execution and funding comparison & & \\[1.8ex]
Dealer/venue intervention with endogenous market response & & \\[1.8ex]
Net value under empirically calibrated disclosure costs & & \\[1.8ex]
Independent implementation and replication & & \\[1.8ex]
\bottomrule
\end{tabular}
\end{center}
\noindent Further data cannot by themselves prove a general identification or
decision-value theorem. Its exact statement, proof and literature priority
require separate review. No publication rating follows from the returned run.

% END INFORMATION DECISION REVISION 20261029

% BEGIN FINITE HORIZON COST 20261029
\section*{Finite-horizon information, cohort ambiguity, and net inventory-service value}

\paragraph{Status and contribution boundary.}
This appendix supplies a restricted analytic separation and its financial
consequence. The proof below is written in full; exact rational and numerical
checks accompany it. It has not received independent proof review and is not
a Lean formalization. Its priority is unconfirmed. Control without unique
identification, task-directed exploration, Gaussian functional estimation,
moment matching, and the value of information are established ideas
\citep{fhc:waarde,fhc:wagenmaker,fhc:donoho,fhc:heinrich,fhc:blackwell}.
The candidate extension is their explicit combination for a finite-horizon
inventory commitment, with a hard stress budget and the information bill
deducted from both capital and profit. It does not close the broader open
theorem of the preceding drafts.

\subsection*{The observation contract and the financial institution}
Fix $H\ge4$, $0\le a<1$, $\kappa>0$ and
$0\le Q\le\overline Q<(1-a)/\kappa$. There are arbitrarily many finite
cohorts, with $q_i\ge0$, $\sum_iq_i=Q$, and $0<\eta_i\le1$:
\[
 x_{t+1}=a x_t+\kappa\bigl(\sum_iq_i y_{i,t}+u_t\bigr),\qquad
 y_{i,t+1}=(1-\eta_i)y_{i,t}+\eta_i x_{t+1}.
\]
Initial states are reset to zero. Let $h_0,\ldots,h_{H-1}$ be the price
response to a unit impulse, with $h_0=\kappa$, and put
$G=\sum_{t<H}h_t$. Positivity implies $h_t\ge0$.
No positive lower bound on $\eta_i$ is imposed.

The observer purchases $N$ independent reset experiments, each lasting $H$
model periods and returning all $H$ taps with independent Gaussian errors of
known variance $\sigma^2$. Thus the tap means are
$\overline h_t=h_t+\sigma Z_t/\sqrt N$, with independent standard normal
$Z_t$. The same observations are available for identifying the cohorts and
for choosing the financial position. One experiment is not one scalar reading:
it returns $H$ readings and costs $c_H$, including the declared operating and
time charges. Total information cost is $F=Nc_H$. This observation contract
is not an ordinary exchange tape; longer horizons cost both time and readings.
Plant parameters, the feedback cap, reset, and stationarity are assumptions.
Measurements are acquired ahead of a common future service-start deadline.
Cash can earn the common reference rate during acquisition; the model does
not include earlier service opportunities. The stated measurement price is
exogenous, not a calibrated self-impact cost for actual market probes.

The institution is a price-taking, unlevered inventory-warehousing desk.
It buys $I\ge0$ units at $P_0$, holds them for $H$ periods, and closes.
It earns a prequoted service fee $pI$ \emph{after} settlement. Round-trip
execution and carry costs are $bI+dI^2$, where $b\ge0$ and $d>0$.
For example, $b=b_{\rm trade}+P_0 r_f H/(365\cdot24\cdot60)$ for
one-minute model periods. The service premium is unavailable as collateral
until after settlement. Both trading costs and $F$ must be prefunded.
Initial cash is $W$, the stressed-settlement loss budget is $B$, and
$|u_t|\le\overline u$, with $\overline u>0$. Assume
$P_0>\overline u\,\overline G$, so the asset price remains positive.
This desk does not alter market feedback, depth, or another participant's
behaviour. The service fee is exogenous, not an equilibrium price.

For the expected-profit objective, future shocks have conditional mean zero
given the observations. Consequently the conditional expected trading price
change is zero. The financial objective, including every declared charge, is
\[
 V(I,F)=g(I)-F,\qquad g(I)=(p-b)I-dI^2.
\]
This is expected service profit with a hard stress constraint, not a Kelly
growth theorem or a return-prediction result. There are no additional training
or validation tapes, no leverage, uncertain fills, or changing inventory.

\begin{lemma}[Sharp prior bounds and exact prefunded capacity]\label{fhc:capacity}
Write
\[
 \underline G=\kappa\sum_{t<H}a^t,\qquad
 \overline G=\kappa\sum_{t<H}(a+\kappa\overline Q)^t.
\]
Then $\underline G\le G\le\overline G$, and both endpoints are attained
in the declared class. For a gain bound $C\ge G$ and
$0\le F\le\min(B,W)$, define
\[
 \begin{gathered}
 R(v,r)=\frac{2v}{r+\sqrt{r^2+4dv}},\qquad I_u=\frac{(p-b)_+}{2d},\\
 I(C,F)=\min\{I_u,R(B-F,\overline u C+b),R(W-F,P_0+b)\}.
 \end{gathered}
\]
This is the optimal expected-profit position subject to the stated stress
reserve and prefunding constraints. The zero-violation no-observation robust optimum is
$I_{\rm nd}=I(\overline G,0)$, not an arbitrarily conservative benchmark.
The response-informed oracle is $I_\circ=I(G,0)$; it knows no future shocks.
\end{lemma}

\begin{proof}
Let $k_j=\sum_iq_i\eta_i(1-\eta_i)^j$. The normalized price response is
the renewal sequence with defective interarrival probabilities
$s_1=a+\kappa k_0$ and $s_j=\kappa k_{j-1}$ for $j\ge2$.
Their total mass is $a+\kappa Q<1$, with the remaining mass placed at
infinity. The cumulative interarrival distribution at integer $j\ge1$ is
\[
 a+\kappa\sum_iq_i[1-(1-\eta_i)^j].
\]
It lies between $a$ and $a+\kappa\overline Q$. Coupling interarrivals by
the same independent uniforms orders every renewal time and hence the
number of renewals before $H$. Taking expectations proves the two bounds.
The lower endpoint is the no-feedback economy; the upper is one cohort
with $Q=\overline Q$ and $\eta=1$. This also proves sharpness of the
no-observation stress bound.

At date $t\le H$, the price change is
$x_t=\sum_{j<t}h_j u_{t-1-j}$, so
$|x_t|\le\overline u\sum_{j<t}h_j\le\overline uG$.
All costs paid by that date are at most the total prefunded costs.
The constant adverse sequence $u_t=-\overline u$ attains the total price
loss $I\overline uG$ at settlement. Before the service fee is credited,
worst loss from initial wealth is therefore exactly
$I\overline uG+bI+dI^2+F$. The declared robust reserve requirement is
$I\overline uC+bI+dI^2+F\le B$; prefunding requires
$P_0I+bI+dI^2+F\le W$. These increasing quadratics yield the two roots.
Finally $g$ increases up to $I_u$ and decreases after it, proving the
optimizer and the two benchmarks. This exact reserve contract is not a
peak-to-trough drawdown or a repeated-deployment guarantee.
\end{proof}

\begin{lemma}[No data at the same error budget]\label{fhc:nodata-chance}
Let $I_f=\min\{I_u,R(W,P_0+b)\}$ and allow a data-free rule to randomize,
with probability of a stress-budget violation at most $\gamma$ in every
economy of the class. Its greatest expected service profit is
\[
 V_{{\rm nd},\gamma}=(1-\gamma)g(I_{\rm nd})+\gamma g(I_f).
\]
It is attained by the displayed two-position lottery. This is a theoretical
benchmark, not a recommended trading policy. The corresponding response-
informed same-error-budget oracle has value
$(1-\gamma)g(I_\circ)+\gamma g(I_f)$.
\end{lemma}

\begin{proof}
Truncating an action at $I_f$ weakly improves profit and lowers risk, so
restrict to $0\le I\le I_f$. At the attainable worst economy with gain
$\overline G$, every $I>I_{\rm nd}$ is unsafe unless $I_{\rm nd}=I_f$,
in which case all larger actions are already dominated. Thus a feasible
lottery puts at most $\gamma$ mass above $I_{\rm nd}$. Below that point
profit is at most $g(I_{\rm nd})$, and above it at most $g(I_f)$.
This gives the upper bound. The two-point lottery attains it and is unsafe
in any economy with probability at most $\gamma$. Replace $\overline G$
by the known $G$ for the oracle statement.
\end{proof}

\begin{theorem}[Uniform cost-aware finite-horizon decision]\label{fhc:decision}
Fix $0<\gamma<1/2$, $z=\Phi^{-1}(1-\gamma)$,
$s_N=\sigma\sqrt{H/N}$, and
\[
 Y=\sum_{t<H}\overline h_t,\qquad
 C_N=\min\{\overline G,\max\{\underline G,Y+zs_N\}\}.
\]
Purchase a fixed, affordable $N$ and take $I_N=I(C_N,F)$.
The following hold uniformly over the declared positive-cohort class,
without dependence on cohort spacing, cohort count, or a memory-rate floor.
\begin{enumerate}
\item With probability at least $1-\gamma$ over measurements, the position
      satisfies the stress budget for \emph{every} bounded future shock path.
\item Put $r_0=\overline u\underline G+b$, $I_{\max}=W/P_0$,
\[
 \begin{aligned}
 L_G&=\frac{\overline u I_{\max}}{r_0}, &
 L_F&=\frac{1}{\min(r_0,P_0+b)},\\
 A&=2z\sigma\sqrt H\,(p-b)_+L_G, & D&=1+(p-b)_+L_F.
 \end{aligned}
\]
With probability at least $1-2\gamma$,
\[
 0\le g(I_\circ)-V(I_N,F)\le A N^{-1/2}+D c_HN.
\]
For sufficiently small $c_H>0$, choosing
$N=\lceil(A/(2Dc_H))^{2/3}\rceil$ is affordable and gives a bound
$3\,2^{-2/3}A^{2/3}(Dc_H)^{1/3}+Dc_H$, when $A>0$.
This is an upper bound from this procedure, not a new minimax optimality claim.
\item Let $I_-=I(G+2zs_N,F)$. A sufficient condition for strictly positive
      expected net information value relative to the strongest no-data rule
      at the \emph{same} error budget is
\[
 \underbrace{(1-2\gamma)g(I_-)-F-V_{{\rm nd},\gamma}}_{\displaystyle\mathcal L_N}>0.
\]
Precisely, $\mathbb E_D[V(I_N,F)]-V_{{\rm nd},\gamma}\ge\mathcal L_N$.
The conclusion can fail at high data costs or a low service premium.
\end{enumerate}
\end{theorem}

\begin{proof}
$Y-G$ is Gaussian with variance $H\sigma^2/N$. On
$\{Y-G\ge-zs_N\}$, which has probability $1-\gamma$, $C_N\ge G$;
clipping preserves this because $G$ lies between the prior bounds.
Lemma~\ref{fhc:capacity} proves simultaneous safety for every future
bounded path on this single event, including the data-selected position.
There is no tail beyond $H$ to estimate, explaining the absence of a
memory-floor assumption. This does not certify an infinite horizon.

On $E=\{|Y-G|\le zs_N\}$, of probability $1-2\gamma$,
$G\le C_N\le G+2zs_N$. The position function is decreasing in each
of $C$ and $F$, so $I_-\le I_N\le I_\circ$ on $E$.
To bound the position difference, implicitly differentiate the stress root:
its derivatives are $-\overline u I/(\overline u C+b+2dI)$ in $C$
and $-1/(\overline u C+b+2dI)$ in $F$.
Clip this root at $I_{\max}$ without changing the final minimum, since
the prefunding root is at most $I_{\max}$. The clipped root is
$L_G$-Lipschitz in $C$ and $1/r_0$-Lipschitz in $F$.
The prefunding root is independent of $C$ and
$1/(P_0+b)$-Lipschitz in $F$. Minimum with $I_u$ preserves these
common Lipschitz bounds. Hence
\[
 0\le I_\circ-I_N\le 2L_Gzs_N+L_FF.
\]
On $[0,I_u]$, $0\le g'\le(p-b)_+$. Multiplication and addition of
the paid fee give the regret bound. Minimizing
$AN^{-1/2}+Dc_HN$ over positive real $N$ gives the stated root and
value. Rounding up adds at most $Dc_H$; affordability follows as
$c_HN=O(c_H^{1/3})\to0$. The case $A=0$ has no positive-gain
information task and does not need this optimization.

On $E$, $g(I_N)\ge g(I_-)$, while $g(I_N)\ge0$ always because
$0\le I_N\le I_u$. Thus
$\mathbb E_D g(I_N)\ge(1-2\gamma)g(I_-)$.
Subtracting $F$ and the same-error-budget no-data value proves the last assertion.
Conditional mean-zero future flows justify the stated expected-profit
interpretation; their distribution is not learned from extra tapes.
\end{proof}

\begin{theorem}[A two-cohort identification obstruction under the same observations]\label{fhc:pair}
Fix $Q>0$, $a+\kappa Q<1$, and $0<\eta_0<1$. For sufficiently small
$\delta>0$, compare the two economies
\[
\begin{array}{c|cc}
 &\text{normalized weights}&\text{adaptation rates}\\ \hline
\mathrm A_\delta &(1/2,1/2)&(\eta_0-\delta,\eta_0+\delta)\\
\mathrm B_\delta &(1/5,4/5)&(\eta_0-2\delta,\eta_0+\delta/2).
\end{array}
\]
Actual weights are $Q$ times those displayed. For fixed $H\ge4$,
$D_H(\delta)=\|h^{\rm B}_{0:H-1}-h^{\rm A}_{0:H-1}\|_2
=\Theta(\delta^3)$. The exact minimax error for distinguishing these
two \emph{known} alternatives from $N$ reset experiments is
\[
 e_N=\Phi\!\left(-\frac{\sqrt N D_H(\delta)}{2\sigma}\right).
\]
Therefore error at most $\beta<1/2$ requires
$N\ge[2\sigma\Phi^{-1}(1-\beta)/D_H(\delta)]^2
=\Theta(\delta^{-6})$.
In particular any recovery of the two normalized weights to sup-norm error
less than $0.15$ with probability at least $1-\beta$ in both worlds must
satisfy this testing requirement.
Both worlds contain exactly two distinct rates. Both have identical infinite
summed response, but their finite-horizon gains are not asserted identical.
\end{theorem}

\begin{proof}
The two rate distributions have equal mass, mean $\eta_0$, and variance
$\delta^2$. Their third centered moments are respectively $0$ and
$-3\delta^3/2$. Thus their kernels $k_j$ agree at $j=0,1$ and
$k^{\rm B}_2-k^{\rm A}_2=-3Q\delta^3/2$. The response recursion
gives equal $h_0,h_1,h_2$ and
\[
 h^{\rm B}_3-h^{\rm A}_3=-3\kappa^2Q\delta^3/2,
\]
which is a lower bound for $D_H$ in absolute value.

For an upper bound put $f_t(\eta)=\eta(1-\eta)^t$.
Its third derivative is
$3t(t-1)(1-\eta)^{t-2}
-\eta t(t-1)(t-2)(1-\eta)^{t-3}$, with absent low-degree terms zero.
The sum over $t\ge0$ of its absolute value is at most $12/\eta^3$.
For $\delta\le\eta_0/4$, all intervening rates are at least
$v=\eta_0/2$. Taylor's integral remainder in $\ell_1$ is therefore at
most $2|\eta-\eta_0|^3/v^3$. The matching first two moments cancel
the quadratic Taylor terms. The sum of the two third absolute centered
moments is $(1+17/10)\delta^3=27\delta^3/10$, giving
\[
 \|k^{\rm B}-k^{\rm A}\|_1\le\frac{27Q}{5v^3}\delta^3.
\]
If $K(z)=\sum k_tz^t$ and $H(z)=\sum h_tz^t$, the recursion gives
$H(z)=\kappa/[1-az-\kappa zK(z)]$. Consequently
$H_{\rm B}-H_{\rm A}=zH_{\rm A}H_{\rm B}(K_{\rm B}-K_{\rm A})$.
The infinite response sums of both economies equal
$S=\kappa/(1-a-\kappa Q)$. Young's inequality yields
$D_H\le S^2\|k^{\rm B}-k^{\rm A}\|_1=O(\delta^3)$.
This proves the two-sided order with constants independent of $\delta$.

For equal-covariance Gaussian laws, the likelihood-ratio test projects the
sample mean onto the difference of the two known means, and thresholds at
their midpoint. The projected normal law gives the displayed error in each
world. Its equal-prior Bayes optimality and symmetry establish minimax
optimality. A weight estimator accurate within $0.15$ distinguishes the
two disjoint normalized-weight balls, proving the identification lower bound.
Finite-horizon gains need not be equal: only their difference is of cubic
order. The test concerns latent cohort recovery, not identification of the
input-output response vector, and does not cover alternative probe designs.
\end{proof}

\begin{corollary}[A priced separation]\label{fhc:priced}
Fix $c_H>0$ and an affordable $N$ for which $\mathcal L_N>0$ at the limiting merged
economy $\delta=0$. For all sufficiently small positive $\delta$, the
financial rule of Theorem~\ref{fhc:decision} has a
strictly positive expected net information value, with its stated safety
probability. Its observation requirement has no dependence on $\delta$.
At this same fixed information price, sufficiently small $\delta$ makes cohort
identification at error $\beta$ unaffordable within the same hard fee budget
$B$: $c_H[2\sigma\Phi^{-1}(1-\beta)/D_H(\delta)]^2>B$.
Any procedure limited to that many affordable observations is covered,
including one that stops early, since it can be viewed as a function of the
full affordable sample. A rational agent may abandon identification and use
the no-data rule. There is no assertion that regulation requires identification.
\end{corollary}

\begin{proof}
Apply the two theorems. The first statement uses the strictly positive
lower bound on net value, rather than an assumption that more information
is automatically beneficial. The second follows from $D_H=\Theta(\delta^3)$.
For a neighbourhood of a fixed $\delta$ where the strict value inequality
holds, continuity preserves it. To claim the inequality for all sufficiently
small $\delta$, it suffices that it hold strictly at the limiting merged
economy: finite $G_H$ and $I_-$ are continuous there. The numerical example
reports both the primary pair and that limiting check. No finite-class grid
is used in the uniform decision guarantee.
\end{proof}

\paragraph{What the result does not establish.}
The direct gain-sum certificate is exactly ordinary Gaussian projection onto
the vector of ones. It is not a new estimator and does not beat that benchmark.
The latent-cohort lower bound is not a lower bound for learning the observable
response vector; comparing it with FIR model-identification results for
$\mathcal H_\infty$ norms would conflate different tasks \citep{fhc:tu}.
No optimal estimation rate over all observation designs is claimed.
The data cost is not an empirically measured exchange probing cost.
An ordinary public-tape bridge, endogenous quotes and depth, incomplete resets,
unknown or misspecified noise, and independent replication remain open.
The information purchase in the example is predeclared. Its positive value
is pointwise at the displayed economy; no prior over unknown economies or
globally optimal Bayesian/robust acquisition policy is claimed.
The fixed horizon, independent measurement contract, and conditional mean-zero
future flows cannot be removed by the simulation. At another horizon the
certificate must be recomputed. A rolling sequence needs simultaneous or
time-uniform coverage; a one-decision guarantee is insufficient.

\paragraph{Literature priority.}
The reviewed sources already establish the general principles used here.
The explicit dollar-cost/capital formulation and the combined result above
are candidate additions to this project, not verified first-in-literature
claims. A specialist must check whether they are an immediate corollary of
existing task-specific or constrained-decision results before publication.

\begin{thebibliography}{9}\small
\bibitem[van Waarde et~al.(2020)]{fhc:waarde} H.~J. van Waarde, J. Eising,
H.~L. Trentelman and M.~K. Camlibel. Data Informativity: A New Perspective on
Data-Driven Analysis and Control. \emph{IEEE TAC} 65, 4753--4768.
\bibitem[Wagenmaker et~al.(2021)]{fhc:wagenmaker} A.~J. Wagenmaker,
M. Simchowitz and K. Jamieson. Task-Optimal Exploration in Linear Dynamical
Systems. \emph{ICML}, PMLR 139, 10641--10652.
\bibitem[Donoho(1994)]{fhc:donoho} D.~L. Donoho. Statistical Estimation and
Optimal Recovery. \emph{Annals of Statistics} 22, 238--270.
\bibitem[Heinrich and Kahn(2018)]{fhc:heinrich} P. Heinrich and J. Kahn.
Strong Identifiability and Optimal Minimax Rates for Finite Mixture Estimation.
\emph{Annals of Statistics} 46, 2844--2870.
\bibitem[Blackwell(1953)]{fhc:blackwell} D. Blackwell. Equivalent Comparisons
of Experiments. \emph{Annals of Mathematical Statistics} 24, 265--272.
\bibitem[Tu et~al.(2018)]{fhc:tu} S. Tu, R. Boczar and B. Recht.
Minimax Lower Bounds for $\mathcal H_\infty$-Norm Estimation.
\emph{arXiv:1809.10855}.
\end{thebibliography}

\section*{Checked synthetic financial illustration}
This illustration uses the frozen primary design and the stronger comparison
recorded in the design amendment. It is exploratory and was developed after
the earlier project runs; it is not independent preregistration or replication.
The two-cohort primary has $a=0.4$, $\kappa=1$, $Q=0.4$, rates
$0.49,0.51$, and equal weights. The known class cap is $\overline Q=0.55$;
$H=60$ one-minute model periods, $N=32$, $\sigma=0.02$, and
$\gamma=0.025$. Initial wealth is 1,000, price 100, stress budget 1,
and flow bound 0.04. The service fee is 0.08 per unit. Execution costs are
$0.002I+0.0002I^2$ and carry is charged at 5 percent annualized.
The assumed all-in reset fee is 0.002, so the learning bill is 0.064.
These are hypothetical model dollars, not an exchange calibration or returns.
The same-error-budget no-data lottery is a theoretical strongest comparator,
not a proposed desk practice. The zero-violation comparisons are also shown.

\begin{center}\footnotesize
\begin{tabular}{p{0.43\textwidth}rrr}
\toprule
Policy & Mean units & Net expected profit & Unsafe reserve (\%)\\
\midrule
No data, zero violations & 1.306 & 0.1007 & 0.00\\
Best no data, same error budget & 1.523 & 0.1171 & 2.50\\
Finite-horizon certificate & 4.553 & 0.2844 & 2.44\\
Standard Gaussian projection & 4.553 & 0.2844 & 2.44\\
Full-response ellipsoid & 4.386 & 0.2718 & 0.00\\
Point estimate (uncertified) & 4.601 & 0.2880 & 50.41\\
Exact gain report, fee 0.05 & 4.669 & 0.3072 & 0.00\\
Response oracle, zero violations & 4.914 & 0.3756 & 0.00\\
Response oracle, same error budget & 5.041 & 0.3851 & 2.50\\
\bottomrule\end{tabular}\end{center}
Profit integrates the future zero-mean flow model and every declared cost;
the data-dependent policy values are also integrated over reading noise by
one-dimensional Gaussian quadrature. Mean units and empirical reserve
violations use 20,000 synthetic reading draws. Oracle and no-data lottery
values are analytic. A reserve violation means some bounded future path can
exceed the budget, not that a sampled market path actually did so.
The 97.5 percent safety statement is over measurement draws for a single
future commitment; it is not certainty, market-wide stability, or time-uniform
coverage. The point rule buys slightly more expected profit by violating the
same risk requirement and is not a valid superiority comparison.
The proposed estimator and standard Gaussian projection coincide exactly.

The net information advantage over the strongest same-error-budget no-data rule is 0.1673. The written theorem's sufficient lower bound is 0.1465, strictly positive even after all costs. At the limiting merged economy the lower bound is 0.1465; continuity therefore supports the small-separation corollary. These bounds are evaluated in floating point with a wide positive margin, not claimed interval-arithmetic certificates.
\begin{center}\small\begin{tabular}{rrr}
\toprule
Separation $\delta$ & Known-pair resets required & Implied fee bill\\
\midrule
0.0800 & 478 & 0.955\\
0.0400 & 3.96e+04 & 79.2\\
0.0200 & 2.84e+06 & 5.67e+03\\
0.0100 & 1.91e+08 & 3.83e+05\\
0.0050 & 1.26e+10 & 2.51e+07\\
0.0025 & 8.15e+11 & 1.63e+09\\
\bottomrule\end{tabular}\end{center}
The known-pair test receives all 60 taps, exactly as the financial rule does.
Response differences are calculated in rational arithmetic before conversion
to floating point. The sample threshold is the Gaussian testing formula, not
a measured runtime or an asserted physical trading-capital requirement.
For the primary $\delta=0.01$, distinguishing the two already-known
alternatives requires about 191 million resets at 5 percent error. Even
this optimistic information fee is far beyond the unit stress budget.
Identifying arbitrary cohorts cannot be easier than this embedded test under
the same fixed query contract. Other experimental designs are not covered.

At fixed $N=32$, the expected net benefit crosses zero at an illustrative reset price of about 0.005825. An exact scalar gain report is worth purchasing only below a fee of about 0.188690 in this example. These are computed break-even prices, not observed reporting costs.
\paragraph{Failures retained.}
The full output includes all 819 cost, horizon and service-fee sensitivity
rows, including negative values and unaffordable information budgets.
A low service premium or expensive measurements can make abstaining from
data acquisition preferable. Independent checks cover 13 algebra, optimizer,
shock-corner, cash-flow and Gaussian-testing properties, including the
same-error-budget no-data optimum checked against a separate linear program.
An initial low-premium LP check was affected by the solver's absolute
objective tolerance; rescaling the objective resolved it, and that failed
log is retained. This is not independent mathematical proof review.

\begin{center}\footnotesize\begin{tabular}{p{0.52\textwidth}rr}
\toprule
Reading/model contract & Gain coverage (\%) & Unsafe reserve (\%)\\
\midrule
Independent Gaussian, valid & 97.36 & 2.63\\
Correlated taps, wrong variance & 71.16 & 28.84\\
Correlated taps, correct variance & 97.36 & 2.63\\
Negative reading bias, invalid & 0.00 & 100.00\\
Future feedback changes, invalid & 0.00 & 100.00\\
\bottomrule\end{tabular}\end{center}
The valid Gaussian row has finite Monte Carlo variation around its nominal
97.5 percent coverage. Using the correct correlated variance restores the
specified measurement contract at the cost of a larger margin. Bias or
future parameter change is not repaired by a larger sample under the wrong
model. A signed kernel such as $(1,-0.9)$ has sum $0.1$ and absolute sum
$1.9$, so positivity is indispensable to using the plain sum. The same
certificate used repeatedly without correction does not retain its
one-decision guarantee.

\paragraph{What changed economically.}
In this model, the desk can choose a larger inventory commitment while
paying for information and retaining the same permitted unsafe-action
probability as its strongest no-data comparator. This is an example of
positive decision value, not evidence that information always pays, that
trading strategies must be identified, or that a real venue's liquidity and
welfare improve. Realistic observation costs and a public-data bridge remain
requirements for an empirical finance claim.

% END FINITE HORIZON COST 20261029

\end{document}
